\documentclass[11pt]{article}

\usepackage[a4paper,margin=2cm]{geometry}
\usepackage{amsmath,amssymb,mathtools}
\usepackage{graphicx}
\usepackage{xcolor}
\usepackage{enumitem}
\usepackage[T1]{fontenc}
\usepackage[utf8]{inputenc}

\graphicspath{{figures/}}
\setlength{\parindent}{0pt}
\setlength{\parskip}{0.45em}
\setlist{nosep,leftmargin=*}

\newcommand{\sourcefig}[2]{%
  \begin{center}
    \includegraphics[width=#2\linewidth]{#1}
  \end{center}}
\newcommand{\dd}{\mathrm{d}}
\newcommand{\Dt}{\frac{\mathrm{D}}{\mathrm{D}t}}
\newcommand{\vect}[1]{\boldsymbol{#1}}
\newcommand{\jump}[1]{\left[#1\right]_{1}^{2}}
\newcommand{\avg}[1]{\left\langle #1\right\rangle}

\begin{document}

\section*{Gasdynamics recap}
\textit{Faithful LaTeX transcription of the 65 numbered source pages.}
The source-page comments are retained throughout so that the transcription can
be checked against the original images.  Extracted drawings are unmodified
pixel crops; their provenance is recorded in \texttt{figures/manifest.json}.

% SOURCE PAGE 1
\section*{Intro}
\textit{giovedì 7 novembre 2024 08:08}

\subsection*{SPREADING OF DISTURBANCES IN COMPRESSIBLE FLOWS}

The flow field is characterized by its velocity \(\vect{V}\) (vectorial
field), its pressure \(p\) (scalar field), and its density/temperature
\(\rho\) and \(T\) (scalar field). Those 3 parameters define the STATE OF THE
FLUID and must be known in order to solve the flow field.

\begin{itemize}
  \item UNIFORM FLOW FIELD: \(\vect{V},p,T\) are constant in time and space;
  \item STEADY FLOW FIELD: \(\vect{V},p,T\) are constant only in time (there
  is a gradient);
  \item UNSTEADY FLOW FIELD:
  \(\vect{V}(\vect{x},t),p(\vect{x},t),T(\vect{x},t)\).
\end{itemize}

We are interested in how changes of the ``primitive variables'' (i.e.
\(\vect{V},p,\rho\)) are propagated. In fact, can be proven that an
infinitesimal disturbance propagate inside the field with a finite speed,
which is the LOCAL SPEED OF SOUND.

\textcolor{blue}{\textbf{DEF WAVE:} Something that propagates with a finite
velocity and changes the state (in our case, the gas).}

There fore in gas, infinitesimal perturbations are carried by waves called
MACH WAVES.

Considering a finite disturbance (such as a piston moving inside a pipe filled
with gas), we can consider every infinitesimal velocity variation as an
infinitesimal source of Mach waves. The resulting finite wave, made up of
infinite Mach waves, can be of two kinds:
\begin{itemize}
  \item COMPRESSION WAVES: \(p\uparrow,T\uparrow\), hence \(a\uparrow\);
  \item EXPANSION WAVES: \(p\downarrow,T\downarrow\), hence \(a\downarrow\).
\end{itemize}
These are continuous isentropic waves because the Mach waves are isentropic as
well.

Considering the compression phenomenon, coalescence can occur:
\(\dd a>0\), hence the next Mach wave is faster and reaches the preceding one.
The compression wave's width shrinks up to a very thin layer (a discontinuity
in our model), forming a SHOCK WAVE (non-isentropic compression).

\textcolor{green}{\textbf{Remainder:}}\quad
\[
a=\sqrt{\frac{\gamma p}{\rho}}=\sqrt{\gamma R T},
\qquad R=\frac{\mathcal{R}}{mm}
\]
for an ideal gas. It is a measure of the thermodynamic state.

\subsection*{Convecting/Material derivative}
Given the tensor field \(\vect{y}(\vect{x},t)\),
\[
\Dt \vect{y}:=\frac{\partial \vect{y}}{\partial t}
 +\vect{u}\mathbin{\cdot}\nabla \vect{y}.
\]
It describes the time-rate change of \(\vect{y}\) associated to a material
element that is subjected to a space--time-dependent velocity field
\(\vect{u}\).

% SOURCE PAGE 2
In fluid dynamics, the velocity field \(\vect{u}\) is the flow velocity field, while
\(\vect{y}\) can be \(p,\rho,T,\ldots\).  In this case the convecting
derivative describes the change of \(\vect{y}\) of a certain fluid element
with time as it follows a trajectory.

% SOURCE PAGE 3
\section*{1-D Unsteady Homentropic Flows}
\textit{7 Nov 2024 09:29}

\subsection*{HOMENTROPIC MODEL}
\textcolor{blue}{\textbf{DEF Entropy}}\quad
\[
s=\frac{p}{\rho^\gamma},\qquad \gamma=\frac{c_p}{c_v}.
\]
If the process is reversible (i.e. compression/expansion waves),
\[
s=\frac{p}{\rho^\gamma}=\mathrm{const}
\qquad\text{(Poisson's Law).}
\]
An isentropic flow has \(s\) constant in time, but its spatial distribution can
vary.  A homentropic flow has \(s\) constant in both time and space.

\subsection*{BALANCE LAWS}
We use a fixed-volume theory (Euler description), with
\(\dd V=A\,\dd x\).  The gas is perfect, non-viscous, and adiabatic.

\sourcefig{page_03_fig_01.png}{0.55}

\textcolor{green}{\textbf{Mass Conservation.}}\quad
The temporal derivative of the total mass in a fixed volume is the net mass
flux through its boundary:
\[
\frac{\dd}{\dd t}\int_{x_1}^{x_2}\rho\,\dd x
=\left.\rho u\right|_{x_1}-\left.\rho u\right|_{x_2}
=-\left[\rho u\right]_{x_1}^{x_2}.
\]
Considering \(\rho u\) continuous, the Green--Gauss theorem applies:
\[
\int_{x_1}^{x_2}\left[\frac{\partial\rho}{\partial t}
+\frac{\partial(\rho u)}{\partial x}\right]\dd x=0.
\]
This gives the divergence form
\[
\rho_t+(\rho u)_x=0,
\]
and the expanded quasi-linear form
\[
\rho_t+u\rho_x+\rho u_x=0.
\]
Equivalently,
\[
\Dt\rho+\rho u_x=0.
\]

\textcolor{blue}{\textbf{Reminder.}}\quad
A quasi-linear equation has coefficients of the highest-order derivatives that
depend at most on derivatives of the previous order.

\textcolor{yellow}{\textbf{Momentum Conservation.}}\quad
Only pressure forces are retained. The total temporal derivative of the linear
momentum inside a volume is equal to the sum between the net momentum flux and
the pressure forces, both at the boundary. The quasi-integral relation begins
on this page and continues on the next page.

% SOURCE PAGE 4
The momentum equation is
\[
\frac{\dd}{\dd t}\int_{x_1}^{x_2}\rho u\,\dd x
=\left[\rho u^2+p\right]_{x_1}
-\left[\rho u^2+p\right]_{x_2}
=-\left[\rho u^2+p\right]_{x_1}^{x_2},
\]
where \(\rho u^2+p\) is the specific momentum flux.  Therefore
\[
(\rho u)_t+(\rho u^2+p)_x=0.
\]
After expansion,
\[
u\bigl[\rho_t+(\rho u)_x\bigr]
+\rho(u_t+u u_x)+p_x=0,
\]
and, using mass conservation,
\[
u_t+u u_x+\frac{1}{\rho}p_x=0,
\qquad
\Dt u=-\frac{1}{\rho}p_x.
\]

\textcolor{blue}{\textbf{Energy Conservation.}}\quad
\[
\frac{\dd}{\dd t}\int_{x_1}^{x_2}\rho e\,\dd x
=-\left[(\rho e)u+p u\right]_{x_1}^{x_2}.
\]
The two terms are the total-energy flux and the instantaneous work of the
pressure forces.  In divergence form,
\[
(\rho e)_t+[(\rho e)u+p u]_x=0.
\]
Here
\[
\rho e=\rho\left(e_i+\frac{u^2}{2}\right)
=\rho\left(c_vT+\frac{u^2}{2}\right)
=\rho\left(\frac{1}{\gamma-1}\frac{p}{\rho}
+\frac{u^2}{2}\right),
\]
with internal and kinetic energy separated.

The expanded equation written in the source is
\[
\begin{aligned}
&\rho_t\left(e_i+\frac{u^2}{2}\right)
+\rho\left[e_{i,t}+\left(\frac{u^2}{2}\right)_t\right]
+\rho_x\left(e_i+\frac{u^2}{2}\right)u+p_xu\\
&\quad+\rho\left[e_{i,x}+\left(\frac{u^2}{2}\right)_x\right]u+p u_x
+\rho\left(e_i+\frac{u^2}{2}\right)u_x=0.
\end{aligned}
\]
Using the chain rule, the mass law, and the momentum law,
\[
\rho\frac{\mathrm{D}e_i}{\mathrm{D}t}+p u_x
+\left(e_i+\frac{u^2}{2}\right)
\left[\frac{\mathrm{D}\rho}{\mathrm{D}t}+\rho u_x\right]
+\rho\left[\frac{\mathrm{D}}{\mathrm{D}t}\left(\frac{u^2}{2}\right)
+\frac{1}{\rho}p_xu\right]=0,
\]
hence
\[
\rho\frac{\mathrm{D}e_i}{\mathrm{D}t}
+p u_x
+\rho u\left[\frac{\mathrm{D}u}{\mathrm{D}t}
+\frac{1}{\rho}p_x\right]=0,
\]
so the quasi-linear energy equation is
\[
\rho\frac{\mathrm{D}e_i}{\mathrm{D}t}+p u_x=0,
\qquad
\rho(e_{i,t}+u e_{i,x})+p u_x=0.
\]
For the homentropic model, Poisson's law can be used instead of the energy
equation.

% SOURCE PAGE 5
\subsection*{Summary}
For a 1-D unsteady homentropic flow, the source summarizes
\[
\begin{aligned}
&\rho_t+u\rho_x+\rho u_x=0,\\
&u_t+u u_x+\frac{1}{\rho}p_x=0,\\
&\frac{p}{\rho^\gamma}=s=\mathrm{const},\\
&\rho(e_{i,t}+u e_{i,x})+p u_x=0.
\end{aligned}
\]
Now we want to rewrite the problem in terms of \(a\), because the equations
become easier to manage. To do so, the sound-speed relation is
\[
a^2=\gamma\frac{p}{\rho}.
\]
Thus
\[
2\ln(a)=\ln(\gamma)+\ln(p)-\ln(\rho),
\]
and logarithmic differentiation gives
\[
2\frac{\dd a}{a}=\frac{\dd p}{p}-\frac{\dd\rho}{\rho}.
\]
Moreover,
\[
\ln(p)-\gamma\ln(\rho)=\ln(s)
\qquad\Longrightarrow\qquad
\frac{\dd s}{s}=\frac{\dd p}{p}-\gamma\frac{\dd\rho}{\rho}.
\]
For a homentropic flow, \(\dd p/p=\gamma\,\dd\rho/\rho\).  Also,
\[
\frac{\dd s}{s}=\frac{1}{p}\left(\dd p-a^2\dd\rho\right),
\qquad
\dd\rho=\frac{\dd p}{p}\frac{\rho}{\gamma}
=\frac{\dd p}{a^2}.
\]
Combining these relations,
\[
\begin{aligned}
2\frac{\dd a}{a}
&=\frac{\dd p}{p}-\frac{\dd p}{a^2\rho}\\
&=\frac{\dd p}{\rho a^2}
\left(\frac{\rho a^2}{p}-1\right)
=\frac{\dd p}{\rho a^2}(\gamma-1).
\end{aligned}
\]
With
\[
\delta=\frac{\gamma-1}{2},
\qquad
\dd a=\frac{1}{\rho a}\frac{\gamma-1}{2}\,\dd p,
\qquad
\dd p=\frac{\rho a}{\delta}\,\dd a,
\]
we have
\[
p_t=\frac{\rho a}{\delta}a_t,\qquad
p_x=\frac{\rho a}{\delta}a_x.
\]
Replacing \(\rho a^2=\gamma p\) in the energy equation gives
\[
p_t+u p_x+\rho a^2u_x=0. \tag{III}
\]
Subtracting \(a^2\) times the mass equation (I) gives
\[
p_t-a^2\rho_t+u(p_x-a^2\rho_x)=0
\qquad\Longrightarrow\qquad
\frac{\mathrm{D}s}{\mathrm{D}t}=0.
\]
Entropy is constant along \(\dd x/\dd t=u\); it is transported by the fluid
along a trajectory.  From \(\mathrm{D}s/\mathrm{D}t=0\), equation (III) can
also be written
\[
\frac{1}{a^2}(p_t+u p_x)+\rho u_x=0
\qquad\Longrightarrow\qquad
p_t+u p_x+a^2\rho u_x=0. \tag{III}
\]
Since (I) and (III) are equivalent, we get rid of (I) and rewrite (III) and
(II) using the \(\dd p\propto\dd a\) relation. This gives the system in
\(a,u\):
\[
\boxed{
\begin{aligned}
a_t+u a_x+\delta a u_x&=0,\\
u_t+u u_x+\frac{a}{\delta}a_x&=0.
\end{aligned}}
\]
The term \((\cdot)_t+u(\cdot)_x\) is tied to a trajectory with velocity \(u\).

% SOURCE PAGE 6
The preceding equations describe the temporal variation of \(a\) and \(u\)
following the velocity flow field:
\[
\frac{\mathrm{D}a}{\mathrm{D}t}=-\delta a u_x,
\qquad
\frac{\mathrm{D}u}{\mathrm{D}t}=-\frac{a}{\delta}a_x.
\]
Since there is a variation, and we are following the velocity of the flow
field, this means that there is something which moves inside the gas with a
velocity \(v\ne u\) that changes the state of the gas itself, hence WAVE.

To find the direction of propagation, consider the trajectory direction in
one dimension, the propagation velocity, and the changed quantity.  Starting
from
\[
\begin{aligned}
a_t+u a_x+\delta a u_x&=0,\\
\delta u_t+\delta u u_x+a a_x&=0,
\end{aligned}
\]
sum and difference give the two characteristic families
\[
\boxed{
\begin{aligned}
(a+\delta u)_t+(u+a)(a+\delta u)_x&=0,\\
(a-\delta u)_t+(u-a)(a-\delta u)_x&=0.
\end{aligned}}
\]
Equivalently, \((a\pm\delta u)_t+(u\pm a)(a\pm\delta u)_x=0\), like a
convecting derivative in which the velocity is replaced by \(u\pm a\).

Two families of waves have been found:
\begin{itemize}
  \item \(u+a\), which propagates with \(u+a\) velocity and changes the
  \(a+\delta u\) quantity;
  \item \(u-a\), which propagates with \(u-a\) velocity and changes the
  \(a-\delta u\) quantity.
\end{itemize}
This is so because the equations say that, if we follow the waves with their
velocities, the respective wave variables are constant. Define
\[
R_1:=a-\delta u,\qquad
R_2:=a+\delta u,
\qquad
\lambda_1=u-a,\quad \lambda_2=u+a.
\]
Then
\[
R_{1,t}+\lambda_1R_{1,x}=0,\qquad
R_{2,t}+\lambda_2R_{2,x}=0.
\]
The Riemann invariants \(R_1,R_2\) are constant along characteristics.
Physically, the state changes when the wave passes.
The source also notes that \(R_1\) and \(R_2\) are related to the state both
thermodynamically and kinematically.

\textcolor{blue}{\textbf{Note.}}\quad
The quasi-linear form does not admit discontinuities.  Characteristic lines
account for the possible finite domain.  For \(t=0\) it is pointless to draw
lines coming from \(t<0\); a boundary condition also makes it unphysical to
draw lines from outside the spatial domain.

% SOURCE PAGE 7
\subsection*{\(x\)-\(t\) plane and geometrical interpretation}

\sourcefig{page_07_fig_01.png}{0.58}

At a point \(P\), the state is defined by \(u_P,a_P,s_P\).  In the \(u\)
direction, \(s=s_P\).  Along the \(u+a\) direction,
\[
a+\delta u=a_P+\delta u_P,
\]
and along the \(u-a\) direction,
\[
a-\delta u=a_P-\delta u_P.
\]
For a homentropic flow, entropy is constant in every direction because
\(\mathrm{D}s/\mathrm{D}t=s_t+u s_x=0\).

\sourcefig{page_07_fig_02.png}{0.58}

Imagine knowing the flow state at \(t=0\) (initial condition) and being able
to construct the characteristic lines starting from \(x_1\) and \(x_2\).
These lines meet at \(P\). Boundary-condition characteristic lines are proper
straight lines if and only if the state is uniform and homogeneous; otherwise
they are curves.

\sourcefig{page_07_fig_03.png}{0.68}

At an intersection,
\[
\begin{aligned}
a+\delta u&=a_1+\delta u_1,\\
a-\delta u&=a_2-\delta u_2,
\end{aligned}
\]
so
\[
a=\frac{a_1+\delta u_1+a_2-\delta u_2}{2},
\qquad
u=\frac{a_1-a_2+\delta u_1+\delta u_2}{2\delta}.
\]
Assuming all characteristic lines can be drawn, we are able to compute the
state of the flow through algebraic computation.

\subsection*{SIMPLE WAVES}
For a simple wave, the characteristic lines of one family are rectilinear.
For the \(u+a\) family,
\[
R_1=a-\delta u=\overline{R}_1\quad\forall(x,t),\qquad
R_2=a+\delta u,
\]
and therefore
\[
a=\frac{R_2+\overline{R}_1}{2},
\qquad
u=\frac{R_2-\overline{R}_1}{2\delta},
\]
\[
\lambda_2=u+a
=\frac{R_2-\overline{R}_1}{2\delta}
+\frac{R_2+\overline{R}_1}{2}.
\]

% SOURCE PAGE 8
Moving along \(\dd x/\dd t=\lambda_2\), \(R_2\) is constant and \(u,a\) are
constant as well.  Thus \(\lambda_2=\dd x/\dd t\) is constant and the line
has constant slope \(u+a\).

\sourcefig{page_08_fig_01.png}{0.50}

For the \(u-a\) family,
\[
R_1=a-\delta u,\qquad
R_2=a+\delta u=\overline{R}_2\quad\forall(x,t),
\]
\[
a=\frac{\overline{R}_2+R_1}{2},
\qquad
u=\frac{\overline{R}_2-R_1}{2\delta}.
\]
Moving along \(\dd x/\dd t=\lambda_1\), \(R_1\), \(u\), and \(a\) are
constant, so \(\lambda_1=\dd x/\dd t\) is constant and the \(u-a\) line is
rectilinear.

\sourcefig{page_08_fig_02.png}{0.62}

\textcolor{blue}{\textbf{Example.}}\quad
For the uniform flow field \(u(x,t)=\widetilde u\),
\(a(x,t)=\widetilde a\), both \(R_1,R_2\) are uniform and every
characteristic line is rectilinear.

% SOURCE PAGE 9
\section*{Centered Expansion}
\textit{8 Nov 2024 12:17}

Now we start to analyze a 1-D homentropic unsteady flow utilizing the shock
tube model. We assume the NON-SLIP condition: the particles in contact with
the piston move with the piston.

\subsection*{Infinite tube, impulsive piston expansion (CENTERED EXPANSION)}
Imagine an infinite tube with a piston. The initial state of the gas is
characterized by \(u_0=0\) and \(a_0\). At \(t=0\) s the piston starts moving
impulsively to the left with velocity \(-v_p\) and it continues to infinity.

\sourcefig{page_09_fig_01.png}{0.60}

We expect the gas on the right side of the piston (the only part of the domain
we are interested in for now) to expand: a region in which \(u\) changes is
formed, then a new equilibrium is established. At \(t=0\) there is a
discontinuity of \(\dot{x}=u=-v_p\) for the particles attached to the piston:
at \(x=0,t=0^-\), \(u=0,a=a_0\), while at \(t=0^+\),
\(u=-v_p,a=a_p\). We avoid the problem by studying
\(t\geq 0+\varepsilon\), where \(\varepsilon\) is very small.

\sourcefig{page_09_fig_02.png}{0.62}
\sourcefig{page_09_fig_03.png}{0.62}

At \(t=0\), the gas has a uniform state: \(a_0\) and \(-a_0\) are the
characteristic lines at any \(x\) (simple waves), in particular also at
\(x=0\), where the piston has not moved. At \(t=0+\varepsilon\) the piston is
moving. At \(x=0\), two finite states coexist because it is a singularity:
\begin{itemize}
  \item \(u=0\), \(a=a_0\): state 0;
  \item \(u=-v_p\), \(a=a_p\): state 1.
\end{itemize}
This perturbation generates two wave fronts, both in the \(u+a\) direction,
since the zone behind the piston is not reachable. The first front has
\[
a+\delta u=a_0,\qquad
\frac{\dd x}{\dd t}=a_0=\lambda_2.
\]
The second front has
\[
a_p+\delta(-v_p),\qquad
\frac{\dd x}{\dd t}=-v_p+a_p=\lambda_2.
\]
The invariant of the \(u-a\) family gives
\[
a_p-\delta u_p=a_p+\delta v_p=a_0,
\qquad
\boxed{a_p=a_0-\delta v_p}.
\]
Notice that \(a_p<a_0\), as expected since this is an expansion
(\(p_1<p_0\), hence \(T_1<T_0\)). Now we can compute the slope of the second
front.
Thus the second front has slope
\[
\frac{\dd x}{\dd t}=-v_p+a_0-\delta v_p
=a_0-(1+\delta)v_p.
\]

% SOURCE PAGE 10
The two fronts delimiting states 0 and 1 are
\[
\left.\frac{\dd x}{\dd t}\right|_0=\lambda_2|_0=a_0,
\qquad
\left.\frac{\dd x}{\dd t}\right|_1
=\lambda_2|_1=a_0-(1+\delta)v_p.
\]
These two lines, both starting from \(x=0\) because of the impulsive law of
the piston, define an ``expansion fan'' in which the state of the gas is
unsteady. Physically speaking, an infinite number of infinitesimal
disturbances are generated at the same spot. For every infinitesimal
disturbance a Mach wave is generated, which infinitesimally changes the state.
To achieve a finite variation, an infinite number of infinitesimal variations
are needed. In this case these infinitesimal variations are expansions.
For \(\dd a<0,\dd u<0\),
\[
\dd(u+a)<0,\qquad \dd(u-a)<0,
\]
so the characteristic lines turn counterclockwise.

\sourcefig{page_10_fig_01.png}{0.62}

The last thing to know is what happens inside the fan. Consider a generic
\((\widetilde x,\widetilde t)\) inside the unsteady region. The simple-wave
relations are
\[
\frac{\widetilde x}{\widetilde t}=u+a,
\qquad
u-\delta a=a_0.
\]
Consequently,
\[
u=\frac{1}{1+\delta}
\left(\frac{\widetilde x}{\widetilde t}-a_0\right),
\qquad
a=a_0+\delta u.
\]
As expected, \(u\) and \(a\) are functions of \(x\) and \(t\). In particular,
\(u\) and \(a\) are linear with respect to \(x\).

\subsection*{Influence of \(-v_p\)}
Considering the \(u+a\) direction,
\[
-v_p+a_p=\lambda_2=a_0-(1+\delta)v_p
\]
is the slope of the tail front.
\begin{enumerate}
\item If \(\lambda_2>0\), then \(v_p<a_0/(1+\delta)\) and
\(v_p<a_p\): the expansion is subsonic, \(M=v_p/a_p<1\).
\item If \(\lambda_2=0\), then \(v_p=a_0/(1+\delta)\): the case is sonic,
\(M=1\).
\item If \(\lambda_2<0\), then \(v_p>a_0/(1+\delta)\) and
\(v_p>a_p\): the case is supersonic, \(M>1\).
\end{enumerate}

\sourcefig{page_10_fig_02.png}{0.96}

% SOURCE PAGE 11
At the limiting case \(a_p=0\), the state in contact with the piston is a
vacuum and the expansion is total:
\[
a_p=a_0-\delta v_p=0,
\qquad
v_p\geq\frac{a_0}{\delta}.
\]
If \(v_p=a_0/\delta\), the tail front is attached to the piston (it is
coincident with \(v_p\)). If \(v_p>a_0/\delta\), a vacuum is created and the
continuous gas model is no longer valid.

\[
\delta=\frac{\gamma-1}{2}\in(0,0.6]
\quad\text{(depending on the gas),}
\qquad
|\Delta a|=\delta|\Delta u|.
\]
Hence \(|\Delta a|<|\Delta u|\) in general.
The non-degenerate case has three zones: the region between the piston and
the tail front, the region ahead of the head front (both zero-family/no-wave
regions), and the region between tail and head, which is a one-family flow
containing one wave.

\sourcefig{page_11_fig_01.png}{0.96}

% SOURCE PAGE 12
\section*{Compression}
\textit{13 Nov 2024 19:26}

\subsection*{COMPRESSION WAVE}
Let the piston position be \(x_p(t)\in C^2\), with
\(\dot x_p=v_p\).  The piston moves to the right, into a uniform state
\(u_0=0,a_0,\gamma\).

\sourcefig{page_12_fig_01.png}{0.60}

The simple wave is in the \(\lambda_2=u+a\) family and is generated by the
left piston.  Across it,
\[
a-\delta u=a_0-\delta u_0,
\qquad
a-a_0=\delta(u-u_0).
\]
Thus \(\dd a=\delta\,\dd u\). The gas is pushed so that \(\dd u>0\), hence
\(\dd a>0\), and
\[
\dd p=\frac{p\gamma}{\delta a}\,\dd a>0.
\]
The propagation directions satisfy \(\dd(u+a)>0\), so that the
\(\lambda_2\)-characteristic lines rotate clockwise when followed from
ahead to behind.

\sourcefig{page_12_fig_02.png}{0.62}

It is clear that the Mach waves will interact somehow. The way of interaction
depends on the distribution of velocity and will be investigated later. In
general, the COALESCENCE phenomenon happens. To understand it, consider the
\((x,u)\) plane of a flow perturbed by a compression wave.

\sourcefig{page_12_fig_03.png}{0.76}

When the wave moves inside the flow, \(u_0\) becomes \(u_1\).  Since
\(\dd(u+a)>0\), the head and tail close in time.  The fixed \(\Delta u=u_1-u_0\)
is achieved over a smaller \(\Delta x\), and \(u_x\to-\infty\).  The Euler
equations cannot be used past this blow-up singularity.  The isentropic
model is valid until the blow-up time \(t_{\mathrm{BU}}\).

% SOURCE PAGE 13
For the \(\lambda_2=u+a\) family, write \(v:=u+a\).  Then
\[
a+\delta u=\mathrm{const}
\quad\text{along}\quad \frac{\dd x}{\dd t}=v.
\]
In the simple wave \(u,a\), and hence \(v\), are constant along
\(\dd x/\dd t=v\), giving the Riemann equation
\[
v_t+v v_x=0.
\]
This is the same equation as inviscid Burgers' equation in the turbulence
context.  The initial-value problem is written
\[
\begin{cases}
v_t+v v_x=0
&\text{along }\dd x/\dd t=v,\quad x\in[a,b],\\
v(t_0,x)=v^0(x)&\text{(initial data).}
\end{cases}
\]

\sourcefig{page_13_fig_01.png}{0.62}

For the characteristic ending at \((\widehat x,\widehat t)\), its foot is
\(\xi\) and
\[
\xi=\widehat x-v^0(\xi)\widehat t,
\qquad
v^0(\xi)=v^0(\widehat x-v^0\widehat t).
\]
We want to verify that \(v(t,x)=v^0(x)=v^0(\xi)\) is the solution. This
expression tells us that the value of \(v\) at any given time can be
reconstructed knowing the initial condition, tracing back along the \(v\)
direction to the \(\xi\) foot. Change variable \(x\to\xi\) in the Riemann
equation:
\[
v_t=v^0_\xi\xi_t=v^0_\xi(-v-v_t t),
\qquad
v_x=v^0_\xi\xi_x=v^0_\xi(1-v_x t).
\]
\[
v^0_\xi(-v-v_t t)+v v^0_\xi(1-v_x t)=0.
\]
Since \(v^0_\xi\ne0\) (there must be a variation, otherwise it is a constant
flow),
\[
-v-v_t t+v-vv_x t=0
\qquad\Longrightarrow\qquad
-t(v_t+vv_x)=0.
\]
The source states that this is true for all \(t\geq0\) for the Riemann
equations.  It then gives
\[
v_x=v^0_\xi(1-v_x t)
=\frac{v^0_\xi}{1+v^0_\xi t}.
\]
The blow-up occurs when \(|v_x|\to\infty\). This phenomenon is called
COALESCENCE, and the denominator vanishes when
\[
1+v^0_\xi t=0,
\qquad
t=-\frac{1}{v^0_\xi}.
\]
If \(v^0_\xi>0\) (expansion), this time is negative and there is no
blow-up for \(t\geq0\).  If \(v^0_\xi<0\) (compression),
\[
t_{\mathrm{BU}}=-\frac{1}{v^0_\xi}
=\frac{1}{|v^0_\xi|}.
\]
Since \(v^0_\xi\) can be a function of \(x\), we must choose the maximum value
of \(|v^0_\xi|\) in order to find the first time at which
\(v_x\to-\infty\). Thus
\[
\boxed{t_{\mathrm{BU}}=\frac{1}{|v^0_\xi|_{\max}}}.
\]

% SOURCE PAGE 14
\subsection*{SIMPLE CENTERED COMPRESSION (linear \(v^0_\xi(x)\))}

\sourcefig{page_14_fig_01.png}{0.54}

Consider the evolution of \(v=u+a\) in time for a linear initial profile.

\sourcefig{page_14_fig_02.png}{0.72}

\sourcefig{page_14_fig_03.png}{0.64}

Since \(v(x)\) is linear, the discontinuity is not only in \(v_x\), but also
in \(v(x)\). This means that at \((x_{\mathrm{BU}},t_{\mathrm{BU}})\) there
are infinite values of \(v\), hence of \(u\) and \(a\). All Mach waves
converge at this point.

Along the three \(\lambda_2\) lines meeting at
\((\widehat x,\widehat t)\),
\[
\widehat a+\delta\widehat u
=a_B+\delta u_B
=a_1+\delta u_1
=a_2+\delta u_2.
\]
This is mathematically valid, but at coalescence it has no physical meaning:
infinitely many Mach lines meet at the single blow-up point, the center of
compression.

For a linear distribution on \(x\in[x_1,x_2]\),
\[
v^0(x)=v_1^0+(v_2^0-v_1^0)
\frac{x_1-x}{x_1-x_2},
\]
\[
v_x^0=-\frac{v_2^0-v_1^0}{x_1-x_2}
=\mathrm{const},
\qquad
t_{\mathrm{BU}}=\frac{1}{|v_x^0|_{\max}}
=\frac{x_1-x_2}{v_2^0-v_1^0}>0.
\]
Geometrically,
\[
v^0(x_2)=v_2^0=\frac{x-x_2}{t},
\qquad
v^0(x_1)=v_1^0=\frac{x-x_1}{t},
\]
so
\[
v_2^0-v_1^0=\frac{-x_2+x_1}{t_{\mathrm{BU}}},
\]
which gives the same time.

\sourcefig{page_14_fig_04.png}{0.60}

\subsection*{Cubic symmetric distribution}
\sourcefig{page_14_fig_05.png}{0.66}

For a cubic symmetric distribution, the Mach-wave coalescence does not
involve all the Mach lines at the same time. The maximum derivative gives
\(t_{\mathrm{BU}}=1/|v_x^0|_{\max}\), attained at \(x=x_0\).

% SOURCE PAGE 15
\sourcefig{page_15_fig_01.png}{0.62}
\sourcefig{page_15_fig_02.png}{0.62}

The exact analytical \(v^0(x)\) determines where coalescence starts.  The
blow-up is not necessarily at the foot \(x_0\); it occurs at
\[
x_{\mathrm{BU}}=x_0+v^0(x_0)t_{\mathrm{BU}}.
\]
At blow-up a shock is generated.  At \(t_{\mathrm{BU}}\) its strength is zero:
\[
M_r=\frac{u-w}{a}
=\frac{u^0(x_0)-u^0(x_{\mathrm{BU}})-a^0(x_0)}{a^0(x_0)}
=-1.
\]
The source labels this as the \(u+a\)-family shock.  As time passes, the shock
is fed by the coalescing waves until it reaches its final strength.

\subsection*{Piston Case (Boundary Condition)}
To find \(t_{\mathrm{BU}}\), one still follows \(u_x\to-\infty\).

\sourcefig{page_15_fig_03.png}{0.72}

The piston law is
\[
x_p(t),\qquad \dot x_p(t)=v_p(t),\qquad \ddot x_p(t),
\qquad t\in[0,t_*],
\]
and the piston then moves at constant speed \(v_p(t_*)\) in the source sketch.
Consider \(\widehat t\in[0,t_{\mathrm{BU}}]\) as a generic time preceding
the blow-up time, and \(t_p<\widehat t\) as a generic time describing the
piston state. Then consider \(\widehat x\) as a generic \(x\) following the
\(u_p+a_p\) direction at \(t_p\). As a consequence,
\(x_p\leq\widehat x\leq a_0t\).
At \(t_p\),
\[
x_p(t_p),\qquad v_p(t_p)=\dot x_p(t_p)=u_p
\]
is the boundary condition.  Along the \(u+a\) characteristic,
\[
a_p(t_p)-\delta u_p=a_0,
\qquad
a_p(t_p)=a_0+\delta v_p(t_p).
\]
For the characteristic through \((\widehat x,\widehat t)\),
\[
\frac{\widehat x-x_p(t_p)}{\widehat t-t_p}
=u+a
=u_p+a_0+\delta v_p
=\dot x_p(t_p)+a_0+\delta\dot x_p(t_p).
\]
Thus
\[
\widehat x(\widehat t)
=\bigl[(1+\delta)\dot x_p(t_p)+a_0\bigr]
(\widehat t-t_p)+x_p(t_p).
\]
The derivative along the boundary parameter is
\[
\widehat u_x
=\frac{\dd\widehat u}{\dd t_p}
\frac{\dd t_p}{\dd\widehat x}
=\frac{\dd\widehat u/\dd t_p}
{\dd\widehat x/\dd t_p}
=\frac{\dd\widehat u}{\dd t_p}
\frac{1}{(1+\delta)\ddot x_p(t_p)(\widehat t-t_p)
-(1+\delta)\dot x_p(t_p)-a_0+\dot x_p(t_p)}.
\]
To find the blow-up time, \(|u_x|\to\infty\), hence
\[
\frac{\dd\widehat x}{\dd t_p}=0
\quad\Longleftrightarrow\quad
(1+\delta)\ddot x_p(t_p)(\widehat t-t_p)
-\delta\dot x_p(t_p)-a_0=0.
\]

% SOURCE PAGE 16
If the piston is initially at rest,
\[
\ddot x_p(0)\neq0,\qquad \dot x_p(0)=0,
\]
the blow-up is on the head front and
\[
(1+\delta)\ddot x_p(0)t_{\mathrm{BU}}-a_0=0,
\qquad
\boxed{t_{\mathrm{BU}}=
\frac{a_0}{(1+\delta)\ddot x_p(0)}}.
\]

\subsection*{Summary on How to deal with Compression Problems}
The isentropic Euler and compatibility equations are valid up to blow-up.
This is calculated imposing \(u\pm a\to-\infty\), depending on the family of
the wave. For initial-condition problems,
\[
t_{\mathrm{BU}}=\frac{1}{|v_x^0|_{\max}},
\qquad
v_x^0=\frac{\partial v^0}{\partial x}.
\]
For a piston boundary condition, if the initial velocity is zero and the
initial acceleration is nonzero,
\[
t_{\mathrm{BU}}=\frac{a_0}{(1+\delta)|\ddot x_p(0)|}.
\]
After blow-up, a moving discontinuity is generated and becomes a shock.  If
the compression is cut off, the shock immediately has maximum strength;
otherwise its strength increases until the compression Mach waves merge.

\subsection*{Centered Compression for Piston Problems}
The source gives the centered-compression piston path
\[
x_p(t)=\widetilde x+
\frac{2a_0}{\gamma-1}
\left[(\widetilde t-t)
-\frac{\gamma+1}{2}\left(
\widetilde t^{\,\frac{\gamma-1}{2}}(\widetilde t-t)\right)\right]^{
\frac{2}{\gamma-1}}.
\]
The stated result is that it generates centered compression with
\(t_{\mathrm{BU}}=\widetilde t\).

% SOURCE PAGE 17
\section*{Discontinuities}
\textit{17 Nov 2024 22:53}

We want to find a way to deal with discontinuities inside the flow because we
have noticed that they arise from compression waves. To do so, we cannot use
the quasi-linear form of the Euler equations. Consider the divergence form:
\[
\begin{aligned}
(\rho)_t+(\rho u)_x&=0,\\
(\rho u)_t+(\rho u^2+p)_x&=0,\\
(\rho e)_t+(\rho e u+p u)_x&=0.
\end{aligned}
\]
With enthalpy \(\rho h=\rho e+p\), the energy equation is
\[
(\rho e)_t+(\rho h u)_x=0.
\]
Here
\[
\rho e=\rho e_i+\rho e_k
=\frac{p}{\gamma-1}+\rho\frac{u^2}{2}.
\]
Let
\[
U=
\begin{bmatrix}\rho\\ \rho u\\ \rho e\end{bmatrix},
\qquad
F(U)=
\begin{bmatrix}
\rho u\\
\rho u^2+p\\
\rho h u
\end{bmatrix}.
\]
The source also writes
\[
U_t+F(U)_x=U_t+(uU+P)_x=0,
\qquad
P=\begin{bmatrix}0\\ p\\ p u\end{bmatrix}.
\]

\sourcefig{page_17_fig_01.png}{0.58}

A discontinuity moves with a constant speed \(w\neq u\).  In the Galilean
reference frame fixed on the discontinuity, \(w\) is constant and
\[
\tau=t,\qquad \xi=x-wt,\qquad v:=u-w.
\]
The coordinate transformation gives, since \(w\) is constant,
\[
U_t=U_\tau\tau_t+U_\xi\xi_t
=U_\tau-U_\xi w=U_\tau-(wU)_\xi,
\qquad
U_x=U_\tau\tau_x+U_\xi\xi_x=U_\xi,
\]
and
\[
F(U)_x=(uU+P)_x=(uU+P)_\xi.
\]
Therefore
\[
U_\tau-(wU)_\xi+(uU)_\xi+P_\xi=0
\qquad\Longrightarrow\qquad
U_\tau+\bigl[(u-w)U+P\bigr]_\xi=0,
\]
or, with \(v:=u-w\),
\[
U_\tau+(vU+P)_\xi=0.
\]

\sourcefig{page_17_fig_02.png}{0.62}

To integrate, consider a control volume (space-time interval) fixed with the
discontinuity:
\[
\int_{\xi_1}^{\xi_2}\left[U_\tau+(vU+P)_\xi\right]\dd\xi=0
\qquad\Longrightarrow\qquad
\frac{\dd}{\dd\tau}\int_{\xi_1}^{\xi_2}U\,\dd\xi
+\int_{\xi_1}^{\xi_2}(vU+P)_\xi\,\dd\xi=0.
\]

% SOURCE PAGE 18
Integrating over a fixed control volume in \(\xi\) gives
\[
\int_{\tau_1}^{\tau_2}
\left[
\frac{\dd}{\dd\tau}\int_{\xi_1}^{\xi_2}U\,\dd\xi
+\int_{\xi_1}^{\xi_2}(vU+P)_\xi\,\dd\xi
\right]\dd\tau=0.
\]
Since the states are uniform and the discontinuity is fixed in the
\(\xi\)-\(\tau\) plane,
\[
\left[\int_{\xi_1}^{\xi_2}U\,\dd\xi\right]_{\tau_1}^{\tau_2}
+\int_{\tau_1}^{\tau_2}\left[vU+P\right]_{\xi_1}^{\xi_2}\dd\tau=0,
\]
and the first term is zero. Since \(\tau_1\) and \(\tau_2\) are generic,
therefore
\[
\left[vU+P\right]_1^2=0.
\]
The jump relations are
\[
\boxed{
\begin{aligned}
(\rho v)_1&=(\rho v)_2,\\
(\rho u v+p)_1&=(\rho u v+p)_2,\\
(\rho e v+p u)_1&=(\rho e v+p u)_2.
\end{aligned}}
\]
The two uniform states separated by the discontinuity are connected by
time-independent relations.  Define the relative mass flux
\[
m:=\rho v.
\]
Mass conservation gives \([m]_1^2=0\).

If \(m=0\), then \(\rho\neq0\), \(v=u-w=0\), so \(w=u\). Momentum gives
\([p]=0\), and energy gives \([pu]=0\).  Writing
\[
[\rho e v+pu]=[me+pu]=m[e]+[pu]=0,
\]
hence \([pu]=0\). We can rewrite
\[
[pu]=p_1u_1-p_2u_2
=\frac{p_1+p_2}{2}(u_1-u_2)
+\frac{u_1+u_2}{2}(p_1-p_2)
=\avg{p}[u]+\avg{u}[p],
\]
we obtain \([u]=0\).  This is a contact discontinuity:
\[
u=\mathrm{const},\qquad p=\mathrm{const},
\qquad [\rho],[s],[T]\ \text{may be nonzero}.
\]

If \(m\neq0\), the discontinuity is a shock wave, an internal moving
boundary.

In general, the direction of propagation and the velocity \(w\) depend on
what caused the discontinuity. A general rule is to think that it must
communicate to the flow that a change occurred. In the previous derivation we
supposed \(w>0\) (towards the right); therefore \(w>u_1,u_2\) (\(u_1,u_2\)
could also be negative while \(w>0\)), otherwise the flow would be faster than
the information.

\subsection*{RANKINE--HUGONIOT RELATIONS}
For a moving shock define the relative Mach numbers
\[
M_{ri}=\frac{v_i}{a_i}=\frac{u_i-w}{a_i},
\qquad
M_{r1}=\frac{u_1-w}{a_1},
\qquad
M_{r2}=\frac{u_2-w}{a_2}.
\]
The source gives
\[
M_{r2}^{\,2}
=\frac{\delta+1/M_{r1}^{\,2}}
{\gamma-\delta/M_{r1}^{\,2}},
\]
\[
\frac{p_2}{p_1}
=1+\frac{2\gamma}{\gamma+1}(M_{r1}^{\,2}-1),
\qquad
\frac{\rho_2}{\rho_1}
=\frac{\gamma+1}{\gamma-1+2/M_{r1}^{\,2}}.
\]
State 1 is the one ahead of the shock wave and state 2 is the one behind.

% SOURCE PAGE 19
\[
\frac{T_2}{T_1}
=\left(\frac{a_2}{a_1}\right)^2
=1+\frac{2(\gamma-1)}{(\gamma+1)^2}
(M_{r1}^{\,2}-1)\left(\gamma+\frac{1}{M_{r1}^{\,2}}\right).
\]

The first consideration is \(\lvert M_{r1}\rvert\geq1\).  If
\(\lvert M_{r1}\rvert<1\), then
\[
\frac{p_2}{p_1}<1,\qquad
\frac{\rho_2}{\rho_1}<1,\qquad M_{r2}^2>1,
\]
which is an expansion. Since \(s:=p/\rho^\gamma\),
\[
\frac{s_2}{s_1}
=\frac{p_2}{p_1}\left(\frac{\rho_1}{\rho_2}\right)^\gamma<1.
\]
This is not possible because the shock is a discontinuity (an isolated system
because it is dimensionless), hence \(\dd s>0\) by the second law of
thermodynamics. Only compression waves are physical. As a consequence,
\[
p_2>p_1,\qquad \rho_2>\rho_1,\qquad u_2>u_1
\]
for the orientation \(w>0\), \(v<0\). Indeed,
\[
\rho_1v_1=\rho_2v_2
\qquad\Longrightarrow\qquad
\frac{\rho_1}{\rho_2}=\frac{v_2}{v_1}<1,
\]
and since \(v:=u-w\), \(v_1<v_2\), hence
\(u_1-w<u_2-w\) and \(u_1<u_2\).

The second consideration is \(M_{r2}^2\leq1\): a physical shock changes the
relative flow from supersonic to subsonic.  As \(M_{r1}^2\to\infty\),
\[
M_{r2}^2\longrightarrow\frac{\gamma-1}{2\gamma},
\qquad
\frac{p_2}{p_1}\longrightarrow+\infty,
\qquad
\frac{\rho_2}{\rho_1}\longrightarrow\frac{\gamma+1}{\gamma-1}.
\]
The pressure rise is unlimited while the Mach-number and density ratios are
limited.  At sufficiently large pressure, dissociation occurs and the ideal
gas model becomes invalid.

\subsection*{Infinite tube, impulsive piston compression (SHOCK WAVE)}
Imagine an infinite tube with a piston. The initial state of the gas is
characterized by \(u_0=0\), \(a_0\), and \(\gamma\). At \(t=0\) s the piston
starts moving impulsively to the right with velocity \(v_p>0\), and it
continues to infinity.

\sourcefig{page_19_fig_01.png}{0.62}

\sourcefig{page_19_fig_02.png}{0.62}

Since the variation of velocity is impulsive, there is a discontinuity at
\(t=0\) s and \(x=0\). Since the gas is pushed by the piston to the right, we
expect an instantaneous compression phenomenon on the \(\lambda_2=u+a\)
family, hence a SHOCK. To solve the flow, we need \(w\), which is unknown. In
this simple case, for which \(u_1=v_p\) because of the boundary condition, it
is possible to find a direct formula.

% SOURCE PAGE 20
For the piston compression,
\[
\frac{\rho_1}{\rho_0}
=\frac{\gamma+1}{\gamma-1+2/M_{r0}^{\,2}},
\qquad
\rho_1v_1=\rho_0(u_1-w)=\rho_0(v_p-w)
=\rho_0v_0=-\rho_0w,
\]
and \(M_{r0}=-w/a_0\).  The shock speed satisfies
\[
w^2-(\delta+1)v_p w-a_0^2=0,
\qquad w>0
\quad\text{(\(u+a\) family).}
\]
An iterative procedure is needed for the general state. With
\(M_{r0}\leq-1\),
\[
\rho_1=\rho_0\frac{\gamma+1}{\gamma-1+2/M_{r0}^{\,2}},
\qquad
p_1=p_0\left[1+\frac{2\gamma}{\gamma+1}
(M_{r0}^{\,2}-1)\right],
\]
\[
M_{r1}
=-\sqrt{\frac{\delta+1/M_{r0}^{\,2}}
{\gamma-\delta/M_{r0}^{\,2}}}.
\]
The minus sign is because \(M_{r0}<0\): in the Galilean frame, the flow
travels from right to left.  Since
\[
M_{r1}=\frac{v_p-w}{a_1},
\qquad
a_1=\frac{v_p-w}{M_{r1}},
\qquad
M_1=\frac{v_p}{a_1},
\]
the piston state is supersonic, sonic, or subsonic according as \(M_1>1\),
\(M_1=1\), or \(M_1<1\).

\sourcefig{page_20_fig_01.png}{0.72}

\sourcefig{page_20_fig_02.png}{0.62}
\sourcefig{page_20_fig_03.png}{0.72}

% SOURCE PAGE 21
If \(v_p\to\infty\), then
\[
M_1=\frac{v_p}{a_1}\longrightarrow\frac{\infty}{\infty},
\]
which can be proved to be limited.

% SOURCE PAGE 22
\section*{Shock Tubes and Riemann Problem}
\textit{21 Nov 2024 10:22}

Now that all the possible types of wave have been introduced, we start to
analyze ``real'' configurations.

\subsection*{Piston \lq\lq shock tube\rq\rq}

\sourcefig{page_22_fig_01.png}{0.62}

Consider an infinite tube with two uniform states:
\[
\text{state 1: }u_1=0,\ p_1=p_0,\ a_1=a_0,\qquad
\text{state 4: }u_4=0,\ p_4=p_0,\ a_4=a_0.
\]
At \(t=0\) s the piston moves with an impulsive law and reaches \(v_p>0\)
instantaneously. The left and right sides do not communicate with each other
because a physical barrier is present in between (the piston). It is possible
to compute the two sides separately.

\subsection*{Left: centered expansion}
Assume \(v_p<a_0/(1+\delta)\), so the expansion is subsonic.  The piston
state is \(u_2=v_p>0\) by the non-slip condition. This information must
propagate toward the left; therefore \(\lambda_1=u-a\), and
\(R_1=a-\delta u\) is the invariant. The source then states that we move along
\(\lambda_2=u+a\), and writes
\[
a_1+\delta u_1=a_2+\delta u_2,
\qquad
a_2=a_0-\delta v_p.
\]
Since \(p/\rho=a^2/\gamma\),
\[
\frac{p_2}{\rho_2}=\frac{a_2^2}{\gamma}
<\frac{a_0^2}{\gamma}=\frac{p_1}{\rho_1},
\]
which is consistent with an expansion. To determine exactly \(\rho_2\) and
\(p_2\), the source gives
\[
\begin{cases}
\displaystyle \frac{p_2}{\rho_2}=\frac{a_2^2}{\gamma},\\
\displaystyle \frac{p_2}{\rho_2^\gamma}
=\frac{p_1}{\rho_1^\gamma}
=p_1\left(\frac{a_0^2}{\gamma p_1}\right)^\gamma.
\end{cases}
\]
Thus state 2 is computed.

\sourcefig{page_22_fig_02.png}{0.66}

\subsection*{Right: shock wave}
Since the gas is pushed by the piston to the right, we expect an instantaneous
compression phenomenon on the \(\lambda_2=u+a\) family, hence a SHOCK. This
specific case has already been solved: the shock speed \(w\) satisfies
\[
w^2-(\delta+1)v_p w-a_4^2=0.
\]
We choose \(w>0\) since this is the \(u+a\) family.

\sourcefig{page_22_fig_03.png}{0.62}

% SOURCE PAGE 23
The shock-speed relation is repeated, choosing \(w>0\) since the shock belongs
to the \(u+a\) family:
\[
w^2-(\delta+1)v_p w-a_4^2=0,
\qquad w>0.
\]
\[
w=\frac{-(\delta+1)v_p+\sqrt{v_p^2(\delta+1)^2+4a_4^2}}{2}.
\]
Then, applying Rankine--Hugoniot, the root and relative Mach number are used to
compute state 3:
\[
M_{r4}=\frac{u_4-w}{a_4}=-\frac{w}{a_4},
\]
\[
\frac{\rho_3}{\rho_4}
=\frac{\gamma+1}{\gamma-1+2/M_{r4}^{\,2}},
\qquad
\frac{p_3}{p_4}
=1+\frac{2\gamma}{\gamma+1}(M_{r4}^{\,2}-1),
\]
\[
M_{r3}
=-\sqrt{\frac{\delta+1/M_{r4}^{\,2}}
{\gamma-\delta/M_{r4}^{\,2}}},
\qquad
a_3=\frac{v_p-w}{M_{r3}}.
\]

\sourcefig{page_23_fig_01.png}{0.68}
\sourcefig{page_23_fig_02.png}{0.68}

The boundary condition of the piston allows us to easily compute the two
unknown states 2 (behind the \(u-a\) expansion) and 3 (behind the \(u+a\)
shock), because it imposes \(u_2=u_3=v_p\). To study other impulsive phenomena
without any physical barrier, we introduce the proper shock-tube problem.

\subsection*{SHOCK TUBE}
\sourcefig{page_23_fig_03.png}{0.62}

The physical shock tube has a driver section, a diaphragm, and a driven
section.  In the driver, state 1 has \(u_1=0,p_1,\gamma\); in the driven
section, state 4 has \(u_4=0,p_4,\gamma\), with \(p_1>p_4\).  If instead
\(p_1=p_4,u_1=u_4\), but \(T_1\neq T_4\), an ideal gas has no flow: there is
no heat transfer and no pressure interaction.

We want to understand what happens when the diaphragm is removed instantly.

At \(t=0\) the pressure jump is released.  The left state pushes to the
right, generating a \(\lambda_2=u+a\) shock; the left-going
\(\lambda_1=u-a\) rarefaction (expansion) is needed. Since the expansion is
isentropic and the shock is not, we also expect an entropy wave that travels
between the other two waves to connect the states behind. Certainly, these

% SOURCE PAGE 24
two states have the same pressure and speed; otherwise another acoustic wave
would be necessary. This is not possible since, at any point of the flowfield,
at most three waves are present, \((\lambda_2,u,\lambda_1)\). A \(u\)-family
contact discontinuity is generated. If the initial conditions are chosen
properly, this last wave can be avoided, but in general it is present.

\sourcefig{page_24_fig_01.png}{0.72}

Now that we know the nature of the waves involved, we have to compute the two
unknown states 2--3 and the velocity of the shock. This kind of problems are
called Riemann problems, and the general solution is presented next.

\subsection*{RIEMANN PROBLEM}
It is the problem of the evolution of a discontinuity inside a flowfield. For
simplicity, we always consider the discontinuity at \(t=0\). This can be done
without loss of generality, because the other way to have discontinuities is
due to compression waves after the blow-up time; hence our Riemann problem
starts from \(t=t_{BU}\).

\sourcefig{page_24_fig_02.png}{0.68}

The solution depends on the waves involved, which are not known a priori,
because the number of unknowns and the number of equations change
accordingly.

The general unknowns are three for state 2, \((u,a,p)\), three for state 3,
\((u,a,p)\), and one for every shock, \(w\): thus there are 6--8 unknowns.
The general equations are two for every isentropic wave,
\[
a\pm\delta u=\mathrm{const},
\qquad
\frac{p_i}{p_j}=\left(\frac{a_i}{a_j}\right)^{2\gamma/(\gamma-1)},
\]
two for the contact discontinuity, \(u_+=u_-\) and \(p_+=p_-\), and three
Rankine--Hugoniot equations for every shock. Thus there are also 6--8
equations.

In any case, the exact solution can be found, but not in closed form because
the resulting algebraic system is nonlinear. The following algorithm can be
used.

First, assume the two waves \((\lambda_1=u-a,\lambda_2=u+a)\) to be
isentropic, together with the contact discontinuity. For states 1--2 in the
\(\lambda_1=u-a\) family, the initial conditions \(u_1,a_1,p_1\) give

\[
\begin{cases}
a_1+\delta u_1=a_2+\delta u_2,\\[2pt]
\displaystyle\frac{p_2}{p_1}
=\left(\frac{a_2}{a_1}\right)^{2\gamma/(\gamma-1)}.
\end{cases}
\]
For states 3--4 in the \(\lambda_2=u+a\) family, the initial conditions
\(u_4,a_4,p_4\) give
\[
a_4-\delta u_4=a_3-\delta u_3.
\]

% SOURCE PAGE 25
The second relation for states 3--4 is
\[
\frac{p_4}{p_3}
=\left(\frac{a_4}{a_3}\right)^{2\gamma/(\gamma-1)}.
\]
Across the contact,
\[
u_2=u_3,\qquad p_2=p_3.
\]
Combining the two Poisson equations and \(p_2=p_3\) gives
\[
a_3=z a_2,
\qquad
z=\left(\frac{p_1}{p_4}\right)^{(\gamma-1)/(2\gamma)}
\frac{a_4}{a_1}.
\]
Here \(z\) is a measure of the isentropicity of states 2--3: if \(z=1\),
they are isentropic. Combining the other two equations gives the initial
velocity guess as a function of the initial conditions:
\[
u_2=u_3=v_0
=\frac{z(a_1+\delta u_1)-(a_4-\delta u_4)}
{\delta(1+z)}.
\]

Second, define the proper nature of the waves, dropping the isentropic
assumption:
\begin{enumerate}
\item For \(\lambda_1=u-a\), if \(u_2<v_0\) it is an expansion; if
\(u_2>v_0\) it is a shock, and the Riemann invariant is replaced by the
Hugoniot relations.
\item For \(\lambda_2=u+a\), if \(u_4>v_0\) it is an expansion; if
\(u_4<v_0\) it is a shock, and the Riemann invariant is replaced by the
Hugoniot relations.
\end{enumerate}

Third, compute \(p_2,p_3\) and, if necessary, iterate. Once we know which
expression to use to compute \(p_2,p_3\)---Poisson for an expansion and
Hugoniot for a shock---we can check whether \(p_2=p_3\). We are assuming
\(u_2=u_3=v_0\), so from the unknown states we must compute only \(p,a\) and,
for shocks, \(w\). For isentropic expansions, Poisson's law and the invariant
equation are enough, while for shocks the three Hugoniot equations determine
\(p,a,w\). If \(p_2=p_3\), then \(v_0=u_2=u_3\) and the problem is over.
If \(p_2\ne p_3\), compute numerically \(v_R\) using Newton--Raphson or the
secant method.

For the secant method, which is easier to compute, another initial guess is
needed. If \(p_2>p_3\), choose \(v_1>v_0\), because the contact discontinuity
should move faster to have \(m=0\); if \(p_2<p_3\), choose \(v_1<v_0\), and
vice versa. Then
\[
v_{k+1}=v_k-\frac{v_k-v_{k-1}}{e_k-e_{k-1}},
\qquad k=1,2,\ldots,\qquad
e_k=\left|1-\frac{p_{2k}}{p_{3k}}\right|.
\]
For every iteration, compute also \(p_2\) and \(p_3\) using the proper
expression. Once the condition \(p_2=p_3\) is reached, the problem is over,
since \(p\) and \(u=v_{k+1}\) are known and \(\rho\) can be calculated.

% SOURCE PAGE 26
\section*{\(p\)-\(u\) plane}
To gain insight into Riemann problems, the subspace of solutions can be
represented on the \(p\)-\(u\) plane. To construct it, start from four basic
configurations that vary both in direction of propagation and in the nature
of the wave. For the \(u-a\) case these are RCN and SCN.

First consider an isentropic rarefaction, subsonic in the drawing.

\sourcefig{page_26_fig_01.png}{0.66}

States 3 and 4 coincide and, for the possible contact,
\[
p_2=p_3=p_4,\qquad u_2=u_3=u_4.
\]
The rarefaction relations are
\[
\begin{cases}
a_2+\delta u_2=a_1+\delta u_1=a_1,\\[2pt]
\displaystyle p_4=p_1
\left(\frac{a_2}{a_1}\right)^{2\gamma/(\gamma-1)}.
\end{cases}
\]
Here \(a_2\ne a_4\) if there is a contact. This solution can be achieved if
we fix \(p_2\) using Poisson's law; otherwise a shock would be present.

\sourcefig{page_26_fig_02.png}{0.66}

Imagine lowering the state-4 pressure, always following Poisson's law, until
vacuum is reached at state \(4v\). This is the largest possible expansion
from state 1, and
\[
u_{4,\mathrm{lim}}
=\frac{a_1+\delta u_1}{\delta}
=\frac{a_1}{\delta}.
\]
The value \(a_1\) also determines the starting slope of the curve. At \(4v\),
the RCN curve is tangent to the \(u\)-axis.

Now consider instead a shock, subsonic in the drawing.

\sourcefig{page_26_fig_03.png}{0.66}

Again, states 3 and 4 coincide and, for the possible contact,
\[
p_2=p_3=p_4,\qquad u_2=u_3=u_4.
\]
The Rankine--Hugoniot relations give
\[
\begin{cases}
\displaystyle p_4=p_1\left[1+\frac{2\gamma}{\gamma+1}
\left(M_{1r}^{2}-1\right)\right],\\[6pt]
\displaystyle \rho_2=\rho_1\left[
\frac{\gamma+1}{\gamma-1+2/M_{1r}^{2}}\right],\\[6pt]
\displaystyle (u_4-w)=(u_1-w)\frac{\rho_1}{\rho_2}.
\end{cases}
\]
Here \(\rho_2\ne\rho_4\) if there is a contact.

\sourcefig{page_26_fig_04.png}{0.66}

% SOURCE PAGE 27
\sourcefig{page_27_fig_01.png}{0.66}

In this case there is no upper limit for the curve because \(p_4\) can be
increased.

For the \(u+a\) case, the two configurations are NCR and NCS. Begin with an
isentropic rarefaction.

\sourcefig{page_27_fig_02.png}{0.66}
\[
\begin{cases}
a_1-\delta u_1=a_4-\delta u_4,
\quad\Longrightarrow\quad a_1=a_4+\delta u_4,\\[2pt]
\displaystyle p_4=p_1
\left(\frac{a_4}{a_1}\right)^{2\gamma/(\gamma-1)}.
\end{cases}
\]
States 1 and 2 coincide and, for the possible contact, \(p_2=p_3\) and
\(u_2=u_3\).

\sourcefig{page_27_fig_03.png}{0.66}

There is no upper limit for \(p_4\).

When considering a shock, care is needed because it can only increase
velocities. To consider a case consistent with \(u_1=0\) in the drawings, we
must choose \(u_4<0\).

\sourcefig{page_27_fig_04.png}{0.66}

The limit is set by \(p_4\to0\). Actually, \(p_4=0\) is not possible because,
if it were so, across the shock
\[
p_4=0\quad\Longrightarrow\quad \rho_3v_3=\rho_4v_4=0,
\]
and the shock would be a contact discontinuity. Therefore, since
\([m]=[\rho v]=0\) and \(m\ne0\), the limit is \(p_4\to0\) and
\(u_4\to-\infty\).

\sourcefig{page_27_fig_05.png}{0.66}

Putting everything together gives the subspace of all possible Riemann
problems with \(u_1=0\) on the left. For a different state 1, the graph can
be translated. All slopes are qualitative and depend on states 1--4.

% SOURCE PAGE 28
\sourcefig{page_28_fig_01.png}{0.84}

The full \(p\)-\(u\) plane contains the regions labeled SCN, SCR, NCR, RCR,
SCS, RCS, RCN, and NCS.  The N.C.R. limit is a double expansion reaching
\(p=0\). This is the limiting case because a vacuum is formed.

\sourcefig{page_28_fig_02.png}{0.76}

The \(p\)-\(u\) plane defines the unique path to achieve equilibrium, even if
interactions between waves or discontinuities occur. It does not give
information about what happens in the interaction region.

\subsection*{How to use \(p\)-\(u\) plane to understand evolution of flow}
Consider, as an example, the classic shock-tube problem with \(p_1>p_4\) and
\(u_1=u_4=0\).

\sourcefig{page_28_fig_03.png}{0.62}

Given the two initial states, connect them using two of the four previous
curves. The only possible combination is an NCR curve starting from state 1
and an NCS curve starting from state 4. This means that the solution is RCS.

\sourcefig{page_28_fig_04.png}{0.74}

This procedure can also be useful to analyze different kinds of problems.

\subsection*{Double shock}

% SOURCE PAGE 29
\sourcefig{page_29_fig_01.png}{0.62}
\sourcefig{page_29_fig_02.png}{0.70}
\sourcefig{page_29_fig_03.png}{0.76}
\sourcefig{page_29_fig_04.png}{0.76}

The two shock paths and their interaction are represented in the tube and in
the \(p\)-\(u\) plane.  After interaction the path is the SCS solution.

% SOURCE PAGE 30
\section*{Interaction of waves}
\textit{gioved\`i 5 dicembre 2024, 11:44}

\subsection*{REFLECTION}
\subsubsection*{Expansion Wave}
Consider a finite piston problem with a wall on the right.

\sourcefig{page_30_fig_01.png}{0.70}

Up to the wall interaction, the centered expansion has \(u_1=0\),
\(u_2=v_p\), \(a_1=a_0\), and
\[
a_1-\delta u_1=a_0=a_2-\delta u_2
=a_2+\delta|v_p|,
\qquad
a_2=a_0-\delta|v_p|.
\]
The interaction with the wall happens in a finite time, not instantaneously,
so an interaction zone is present. The wall imposes \(u_3=0\) as a boundary
condition. We expect the velocity to turn clockwise, since before it turned
counterclockwise. The reflected wave is still an expansion wave, but on the
\(\lambda_1=u-a\) family:
\[
a_2+\delta u_2=a_2-\delta|v_p|=a_3.
\]
The source then writes
\[
a_3=a_0+\delta|v_p|.
\]

Instead, now consider a vessel that imposes \(p=\mathrm{const}\) as a
boundary condition.

\sourcefig{page_30_fig_02.png}{0.72}

The process is still isentropic; hence, if \(p=\mathrm{const}\), then
\(\rho\) and \(T\) are constant too, so \(a_3=a_0\). The vessel must
compensate the pressure decrease given by the \(u+a\) expansion wave, so a
compression wave is generated. Along the \(\lambda_2=u+a\) direction,
\[
a_2+\delta u_2=a_2-\delta|v_p|
=a_3+\delta u_3=a_0+\delta u_3,
\qquad
u_3=\frac{a_2-\delta|v_p|-a_0}{\delta}
=2v_p<0.
\]

\subsubsection*{Shock Wave}

% SOURCE PAGE 31
\sourcefig{page_31_fig_01.png}{0.74}

At the wall \(u_3=0\).  The instantaneous interaction produces a
\(\lambda_1=u-a\) shock, which is non-isentropic.  State 1 is known,
\(a_1=a_0,u_1=0\), and state 2 has \(u_2=v_p\).  For an exact impulsive
piston or an iterative calculation, the source uses
\[
a_2-\delta v_p=a_0,
\qquad a_2=a_0+\delta v_p,
\]
and the first guess
\[
w_0=\frac{(u_2+a_2)+(u_1+a_1)}{2}
 =\frac{(v_p+a_0+\delta v_p)+a_0}{2}.
\]
Then
\[
M_{r1}=\frac{u_1-w_0}{a_1}=-\frac{w_0}{a_0},
\qquad
M_{r2}=\sqrt{\frac{\delta+1/M_{r1}^{2}}
{\gamma-\delta/M_{r1}^{2}}}
=\frac{u_2-w_0}{a_2},
\qquad u_2\simeq v_p.
\]
If \(u_{2,1}>v_p\), choose \(w_1<w_0\); otherwise choose \(w_1>w_0\).
Repeat until
\[
|u_{2k}-v_p|<\varepsilon.
\]
Once \(w\) is known, \(M_{r1}\) is known and state 2 follows from the
Rankine--Hugoniot relations.

To compute the reflected shock, apply a similar procedure knowing that
\(u_3=0\); the iterative method is necessary. The resulting \(w\) is lower,
so the phenomenon is not linear.

The manuscript refers to p.~218 of ``PDF Sabbia [ITALIAN]'' for easier
formulas. If a vessel imposes \(p_3=\mathrm{const}\), it can be proved
(p.~120) that
\[
a_3=a_0 e^{s_2-s_1},
\qquad s=\frac{p}{\rho^\gamma}.
\]

\subsection*{Interaction between waves of the same kind}
If two waves of the same kind but opposite family interact, the phenomenon
is the same as reflection against a wall. Expansion waves of the same family
cannot interact because the tail-front velocity of the leading one is the
same as the head-front velocity of the following one. Shock waves of the same
family do interact, generating a Riemann problem.

\sourcefig{page_31_fig_02.png}{0.72}

% SOURCE PAGE 32
\sourcefig{page_32_fig_01.png}{0.74}

Regions I--IV are the two starting regions.

\subsection*{Friedrichs Problems}
Consider an infinite pipe with a piston. We investigate the interaction
between expansions and shocks generated by an impulsive motion law of the
piston itself.

\subsubsection*{Direct problem}
The piston is moved with \(v_p\) and then stopped.

\sourcefig{page_32_fig_02.png}{0.68}

The first acceleration generates a shock; the stop then generates an
expansion wave, bringing the gas the information that the piston has stopped.
Both waves belong to the \(\lambda_2=u+a\) family.

\sourcefig{page_32_fig_03.png}{0.68}

The interaction is not instantaneous. Therefore an entropy layer, not a
contact discontinuity, is generated, and there is an interaction region in
which the characteristic lines are curved. Neglecting all of this, the
problem can be solved as a ``pseudo RP'': the nature of the generated waves
and the states they form can be investigated using the \(p\)-\(u\) plane.

\sourcefig{page_32_fig_04.png}{0.76}

Once states A, B, and C have been computed, it is clear that other acoustic
waves must be generated to establish equilibrium between A and C. Treat A
and C as two initial states of a Riemann problem and use the \(p\)-\(u\) plane
to find the proper waves.

The boundary-condition note is that, for an expansion process and a given
\(\Delta u\),
\[
\Delta p_{\mathrm{exp}}<\Delta p_{\mathrm{shock}},
\]
because a smaller amount of internal energy is needed.

% SOURCE PAGE 33
\sourcefig{page_33_fig_01.png}{0.66}
\sourcefig{page_33_fig_02.png}{0.66}

The RCS solution consists of an expansion in the \(u-a\) family, an entropy
layer/contact on \(u\), and a shock in the \(u+a\) family.

\subsubsection*{Reverse}
\sourcefig{page_33_fig_03.png}{0.66}

The reasoning adopted is exactly the same.

\sourcefig{page_33_fig_05.png}{0.76}

For a compression process and a given \(\Delta u\),
\[
\Delta p_{\mathrm{comp}}>\Delta p_{\mathrm{shock}},
\]
because only internal energy is used.

\sourcefig{page_33_fig_04.png}{0.66}

The SCR solution consists of a shock on the \(u-a\) family, an entropy
layer/contact on \(u\), and an expansion on the \(u+a\) family.

\subsection*{Another example}
\sourcefig{page_33_fig_06.png}{0.64}

Imagine a shock tube closed on the left. The diaphragm is opened, and the
states satisfy
\[
p_1>p_4,\qquad u_1=u_4=0,\qquad a_1=a_4.
\]
The resulting Riemann pattern continues on the next page.

% SOURCE PAGE 34
\sourcefig{page_34_fig_01.png}{0.68}

We expect a Riemann problem with the qualitative solution RNS. No contact
discontinuity is formed because \(a_1=a_4\), hence \(s_1=s_4\).

\sourcefig{page_34_fig_02.png}{0.76}
\sourcefig{page_34_fig_03.png}{0.76}

After a time
\[
t^*=\frac{L}{a_2},
\]
the head front of the \(u-a\) expansion wave reaches the left boundary and a
reflection occurs. The reflected wave is going to interact with the shock,
just as in the piston case; this time the starting states are 5 and 4, both
with \(u=0\).

State 5 is reached by an isentropic wave: the invariants give a \(u+a\)
expansion from 2 to 5. Thus \(R_1=a-\delta u\) over
\(\lambda_1=u-a\) gives
\[
a_2-\delta u_2=a_5,
\qquad
p_5=p_2\left(\frac{a_5}{a_2}\right)^{2\gamma/(\gamma-1)},
\]
where \(u_2\) is known from the Riemann problem.

% SOURCE PAGE 35
\section*{Method of Characteristics}
\textit{mercoled\`i 15 gennaio 2025, 15:53}

We now generalize the problem from a mathematical point of view in order to
construct a numerical method.

\subsection*{Wave Solution of Euler Equations}
Start from the quasi-linear form of the balance laws:
\[
\begin{aligned}
\rho_t+u\rho_x+\rho u_x&=0,\\
u_t+u u_x+\frac{1}{\rho}p_x&=0,\\
p_t+u p_x+\gamma p u_x&=0.
\end{aligned}
\]
Define the primitive-variables vector
\[
U=\begin{bmatrix}\rho\\u\\p\end{bmatrix},
\qquad
U_t+A U_x=0,
\qquad
A=
\begin{bmatrix}
u&\rho&0\\
0&u&1/\rho\\
0&\gamma p&u
\end{bmatrix},
\]
where \(A\) is the Jacobian. Since these equations can be rewritten in terms
of wave variables conserved along characteristic directions, use them to
diagonalize the problem. This operation can be seen as a change of reference
frame.

\sourcefig{page_35_fig_01.png}{0.58}

For the \(k\)-th rotated frame, with \(\theta^k\ne\pi/2\),
\[
U_t
=U_{\xi^k}\xi_t^k+U_{\tau^k}\tau_t^k
=\sin\theta^k U_{\xi^k}+\cos\theta^k U_{\tau^k},
\qquad
U_x
=U_{\xi^k}\xi_x^k+U_{\tau^k}\tau_x^k
=\cos\theta^k U_{\xi^k}-\sin\theta^k U_{\tau^k}.
\]
Hence
\[
(\cos\theta^k I-\sin\theta^k A)U_{\tau^k}
+(\sin\theta^k I+\cos\theta^k A)U_{\xi^k}=0.
\]
Define
\[
\lambda^k:=\frac{\sin\theta^k}{\cos\theta^k}=\tan\theta^k.
\]
If \(\lambda^k\) is a direction of propagation, then \(U_{\tau^k}=0\), so
\[
(-\lambda^kI+A)U_{\xi^k}=0.
\]
We seek a nontrivial solution, giving the eigenproblem
\[
\det(-\lambda^k I+A)=0,
\qquad
\det\begin{bmatrix}
u-\lambda^k&\rho&0\\
0&u-\lambda^k&1/\rho\\
0&\gamma p&u-\lambda^k
\end{bmatrix}=0.
\]
Thus
\[
(u-\lambda^k)\bigl[(u-\lambda^k)^2-\gamma p/\rho\bigr]=0,
\]
so
\[
\lambda_1=u-a,\qquad \lambda_2=u,\qquad \lambda_3=u+a.
\]

% SOURCE PAGE 36
Computing the eigenvectors gives the left row vectors
\[
\ell_1=\begin{bmatrix}0&-\rho a&1\end{bmatrix},
\qquad
\ell_2=\begin{bmatrix}-a^2&0&1\end{bmatrix},
\qquad
\ell_3=\begin{bmatrix}0&\rho a&1\end{bmatrix}.
\]
In particular, we can verify that these are waves:
\[
\begin{aligned}
\ell_1\mathbin{\cdot}(U_t+A U_x)
&=(p_t-\rho a u_t)+(u-a)(p_x-\rho a u_x)\\
&=\dd p-\rho a\,\dd u=\dd\sigma_1=0
\quad\text{along }\lambda_1=u-a,\\
\dd\sigma_2
&=(p_t+a^2\rho_t)+u(p_x+a^2\rho_x)
=s_t+u s_x=0
\quad\text{along }\lambda_2=u,\\
\ell_3\mathbin{\cdot}(U_t+A U_x)
&=(p_t+\rho a u_t)+(u+a)(p_x+\rho a u_x)=0
\quad\text{along }\lambda_3=u+a.
\end{aligned}
\]

Define matrices that will be useful later:
\[
L=
\begin{bmatrix}
0&-\rho a&1\\
-a^2&0&1\\
0&\rho a&1
\end{bmatrix},
\qquad
R=L^{-1}
=\frac{1}{2a^2}
\begin{bmatrix}
1&-2&1\\
-a/\rho&0&a/\rho\\
a^2&0&a^2
\end{bmatrix},
\]
and
\[
\Lambda=LAR=\operatorname{diag}(u-a,u,u+a).
\]

\subsection*{Linear Numerical Approximation (Cauchy--Kowalewskaya)}
The source writes
\[
U_t+L\Lambda R U_x=0
\]
and notes that there are no analytical solutions except for simple waves,
whose characteristic lines are straight.  The linearized isentropic form is
used to build a linearized solution, with no discontinuities.

One possible approach is the Cauchy--Kowalevskaya method, which consists in
approximating the partial derivatives with a finite-difference form.  The
source says that, for reasons known only to Favini, this process is presented
not in the matrix form of the Euler equations but in the conserved-variable
(wave) form.  Perhaps this is because characteristic lines are used to
discretize the domain; the source ends the aside with
\lq\lq Doesn't Matter\rq\rq.

% SOURCE PAGE 37
\sourcefig{page_37_fig_01.png}{0.68}

The first step is to discretize the space--time domain.  With the propagated
variables, fix a constant value
\[
\Delta x=x_j-x_{j-1}
\]
at every point of the domain. The same is not true for \(\Delta t\), which
will depend on a parameter introduced later. Consider the Euler equations in
the form of propagated (conserved) variables:
\[
L(U_t+R\Lambda L U_x)=0
\qquad\Longrightarrow\qquad
L U_t+\Lambda L U_x=0.
\]
The wave equations are
\[
\begin{aligned}
\dd p-\rho a\,\dd u&=0&&\text{along }\lambda_1=u-a,\\
\dd p-a^2\dd\rho&=0&&\text{along }\lambda_2=u,\\
\dd p+\rho a\,\dd u&=0&&\text{along }\lambda_3=u+a.
\end{aligned}
\]
If considered row by row, \(\ell_k U_t+\ell_k A U_x=0\) is the same set
analyzed before.

\sourcefig{page_37_fig_02.png}{0.70}
\sourcefig{page_37_fig_03.png}{0.70}

At \((x_j,t^n)\), linearize with
\(\widehat A=A(x_j,t^n)\), making the characteristics straight:
\[
\Delta p-\widehat\rho\,\widehat a\,\Delta u=0,\qquad
\Delta p-\widehat a^2\Delta\rho=0,\qquad
\Delta p+\widehat\rho\,\widehat a\,\Delta u=0.
\]
Choose
\[
\Delta t^n=\alpha\frac{\Delta x}{\lambda_{j,\max}},
\qquad
\lambda_{j,\max}=\max_j(|u_j^n|+a_j^n),
\qquad 0<\alpha\leq1.
\]
The time step changes after each iteration.  The boundary-condition formulae
shown are valid when the feet of the characteristics lie as in the source
sketch (positive subsonic flow, \(u_j^n>0\), \(u_j^n<a_j^n\)); otherwise the
equations must be modified.

The next step is to use the equations in discrete form; for example, consider
the \(u+a\) direction.

% SOURCE PAGE 38
For the \(u+a\) characteristic,
\[
\Delta p+\widehat\rho\,\widehat a\,\Delta u=0,
\]
which at the grid point is
\[
p(x_j,t^{n+1})-p(x_3,t^n)
+\widehat\rho\,\widehat a
\bigl[u(x_j,t^{n+1})-u(x_3,t^n)\bigr]=0.
\]
Interpolating to \(x_3\) gives
\[
p(x,t^n)=p_j^n-
\frac{p_j^n-p_{j-1}^n}{\Delta x}(x-x_j),
\qquad
p(x_3,t^n)
=p_j^n-\frac{p_j^n-p_{j-1}^n}{\Delta x}
(u_j^n+a_j^n)\Delta t,
\]
\[
u(x,t^n)=u_j^n-
\frac{u_j^n-u_{j-1}^n}{\Delta x}(x-x_j),
\qquad
u(x_3,t^n)
=u_j^n-\frac{u_j^n-u_{j-1}^n}{\Delta x}
(u_j^n+a_j^n)\Delta t.
\]
Substitution in the characteristic equation gives
\[
p_j^{n+1}-\left[p_j^n-
\frac{p_j^n-p_{j-1}^n}{\Delta x}(u_j^n+a_j^n)\Delta t\right]
+\rho_j^n a_j^n\left[u_j^{n+1}-\left(u_j^n-
\frac{u_j^n-u_{j-1}^n}{\Delta x}(u_j^n+a_j^n)\Delta t\right)\right]=0.
\]
The highlighted finite-difference relation is
\[
\frac{p_j^{n+1}-p_j^n}{\Delta t}
+\rho_j^n a_j^n\frac{u_j^{n+1}-u_j^n}{\Delta t}
+(u_j^n+a_j^n)
\left[
\frac{p_j^n-p_{j-1}^n}{\Delta x}
+\rho_j^n a_j^n\frac{u_j^n-u_{j-1}^n}{\Delta x}
\right]=0.
\]

For the \(u-a\) characteristic,
\[
\Delta p-\widehat\rho\,\widehat a\,\Delta u=0,
\]
and the interpolation to \(x_2\) uses
\((u_j^n-a_j^n)\Delta t\):
\[
p(x,t^n)=p_j^n-
\frac{p_{j+1}^n-p_j^n}{\Delta x}(x-x_j),
\qquad
p(x_2,t^n)=p_j^n-
\frac{p_{j+1}^n-p_j^n}{\Delta x}(u_j^n-a_j^n)\Delta t,
\]
\[
u(x,t^n)=u_j^n-
\frac{u_{j+1}^n-u_j^n}{\Delta x}(x-x_j),
\qquad
u(x_2,t^n)=u_j^n-
\frac{u_{j+1}^n-u_j^n}{\Delta x}(u_j^n-a_j^n)\Delta t.
\]
Thus
\[
p_j^{n+1}-\left[p_j^n-
\frac{p_{j+1}^n-p_j^n}{\Delta x}(u_j^n-a_j^n)\Delta t\right]
-\rho_j^n a_j^n\left[u_j^{n+1}-\left(u_j^n-
\frac{u_{j+1}^n-u_j^n}{\Delta x}(u_j^n-a_j^n)\Delta t\right)\right]=0,
\]
and hence
\[
\frac{p_j^{n+1}-p_j^n}{\Delta t}
-\rho_j^n a_j^n\frac{u_j^{n+1}-u_j^n}{\Delta t}
+(u_j^n-a_j^n)
\left[
\frac{p_{j+1}^n-p_j^n}{\Delta x}
-\rho_j^n a_j^n\frac{u_{j+1}^n-u_j^n}{\Delta x}
\right]=0.
\]

Along the \(u\) characteristic,
\[
\Delta p-\widehat a^2\Delta\rho=0,
\]
so
\[
p(x_j,t^{n+1})-p(x_2,t^n)
-\widehat a^2\bigl[\rho(x_j,t^{n+1})-\rho(x_2,t^n)\bigr]=0.
\]
The source begins the interpolation at the foot \(x_2\):
\[
p(x,t^n)=p_j^n-\frac{p_j^n-p_{j-1}^n}{\Delta x}(x-x_j),
\qquad
p(x_2,t^n)=p_j^n-\frac{p_j^n-p_{j-1}^n}{\Delta x}u_j^n\Delta t,
\]
\[
\rho(x,t^n)=\rho_j^n-\frac{\rho_j^n-\rho_{j-1}^n}{\Delta x}(x-x_j).
\]

% SOURCE PAGE 39
Continuing the interpolation,
\[
\rho(x_2,t^n)=\rho_j^n
-\frac{\rho_j^n-\rho_{j-1}^n}{\Delta x}u_j^n\Delta t.
\]
Substitution gives
\[
p_j^{n+1}-\left[p_j^n-
\frac{u_j^n(p_j^n-p_{j-1}^n)}{\Delta x}\Delta t\right]
-(a_j^n)^2\left[\rho_j^{n+1}-\left(\rho_j^n-
\frac{u_j^n(\rho_j^n-\rho_{j-1}^n)}{\Delta x}\Delta t\right)\right]=0,
\]
and therefore
\[
\frac{p_j^{n+1}-p_j^n}{\Delta t}
-(a_j^n)^2\frac{\rho_j^{n+1}-\rho_j^n}{\Delta t}
+u_j^n\left[
\frac{p_j^n-p_{j-1}^n}{\Delta x}
-(a_j^n)^2\frac{\rho_j^n-\rho_{j-1}^n}{\Delta x}
\right]=0.
\]
Putting everything together, the discretized wave-form system is
\[
\begin{cases}
\displaystyle
\frac{p_j^{n+1}-p_j^n}{\Delta t}
-\rho_j^n a_j^n\frac{u_j^{n+1}-u_j^n}{\Delta t}
+(u_j^n-a_j^n)\left[
\frac{p_{j+1}^n-p_j^n}{\Delta x}
-\rho_j^n a_j^n\frac{u_{j+1}^n-u_j^n}{\Delta x}
\right]=0,\\[10pt]
\displaystyle
\frac{p_j^{n+1}-p_j^n}{\Delta t}
-(a_j^n)^2\frac{\rho_j^{n+1}-\rho_j^n}{\Delta t}
+u_j^n\left[
\frac{p_j^n-p_{j-1}^n}{\Delta x}
-(a_j^n)^2\frac{\rho_j^n-\rho_{j-1}^n}{\Delta x}
\right]=0,\\[10pt]
\displaystyle
\frac{p_j^{n+1}-p_j^n}{\Delta t}
+\rho_j^n a_j^n\frac{u_j^{n+1}-u_j^n}{\Delta t}
+(u_j^n+a_j^n)\left[
\frac{p_j^n-p_{j-1}^n}{\Delta x}
+\rho_j^n a_j^n\frac{u_j^n-u_{j-1}^n}{\Delta x}
\right]=0.
\end{cases}
\]
This is the finite-difference approximation of the wave form of the Euler
equations. Using the same approach, the matrix form can also be approximated.
In matrix form, distinguish what comes from the left and from the right for
two reasons: the relative points with respect to \(x_j\) are different, and
at a boundary the boundary condition must be imposed from the right
direction.

Separate characteristics according to direction:
\[
\Lambda^+=\frac{\Lambda+|\Lambda|}{2},
\qquad
\Lambda^-=\frac{\Lambda-|\Lambda|}{2}.
\]
The source writes
\[
U_t+R\Lambda L U_x=0
\quad\Longrightarrow\quad
U_t+R\Lambda^+L\,U_x\big|_{\text{from left}}
+R\Lambda^-L\,U_x\big|_{\text{from right}}=0
\]

% SOURCE PAGE 40
and the upwind finite-difference system
\[
\frac{U_j^{n+1}-U_j^n}{\Delta t^n}
+\frac{1}{\Delta x}\left[
R_j^n\Lambda_j^{+n}L_j^n(U_j^n-U_{j-1}^n)
+R_j^n\Lambda_j^{-n}L_j^n(U_{j+1}^n-U_j^n)
\right]=0.
\]

The final update is
\[
\boxed{
U_j^{n+1}=U_j^n-\frac{\Delta t^n}{\Delta x}\left[
R_j^n\Lambda_j^{+n}L_j^n(U_j^n-U_{j-1}^n)
+R_j^n\Lambda_j^{-n}L_j^n(U_{j+1}^n-U_j^n)
\right].}
\]

% SOURCE PAGE 41
\section*{Godunov Method}
\textit{gioved\`i 16 gennaio 2025, 12:08}

The Godunov method is a numerical approximation that allows
discontinuities.  Starting from
\[
\begin{aligned}
\rho_t+(\rho u)_x&=0,\\
(\rho u)_t+(\rho u u+p)_x&=0,\\
(\rho e)_t+(\rho h u)_x&=0,
\end{aligned}
\qquad
U_t+F(U)_x=0,
\]
with
\[
U=\begin{bmatrix}\rho\\ \rho u\\ \rho e\end{bmatrix},
\qquad
F(U)=\begin{bmatrix}\rho u\\ \rho u^2+p\\ \rho h u\end{bmatrix},
\]
the integral form over a fixed interval is
\[
\frac{\dd}{\dd t}\int_{x_1}^{x_2}U\,\dd x
+[F(U)]_{x_1}^{x_2}=0.
\]

\sourcefig{page_41_fig_01.png}{0.70}

Use cell interfaces in the discretization:
\[
\frac{\dd}{\dd t}\int_{x_{j-1/2}}^{x_{j+1/2}}U\,\dd x
+[F(U)]_{x_{j-1/2}}^{x_{j+1/2}}=0.
\]
Integrating from \(t^n\) to \(t^{n+1}\),
\[
\int_{t^n}^{t^{n+1}}
\frac{\dd}{\dd t}\left(
\int_{x_{j-1/2}}^{x_{j+1/2}}U\,\dd x
\right)\dd t
+\int_{t^n}^{t^{n+1}}
[F(U)]_{x_{j-1/2}}^{x_{j+1/2}}\,\dd t=0,
\]
and therefore
\[
\int_{x_{j-1/2}}^{x_{j+1/2}}
(U^{n+1}-U^n)\,\dd x
+\int_{t^n}^{t^{n+1}}
[F(U)]_{x_{j-1/2}}^{x_{j+1/2}}\,\dd t=0.
\]
To compute the first integral, the mean-value theorem gives an exact solution
up to second order:
\[
\int_{x_{j-1/2}}^{x_{j+1/2}}
(U^{n+1}-U^n)\,\dd x
=\Delta x\left(\overline U_j^{\,n+1}
-\overline U_j^{\,n}\right).
\]

\sourcefig{page_41_fig_02.png}{0.42}

% SOURCE PAGE 42
\[
\overline U_j^{\,n+1}-\overline U_j^{\,n}
+\frac{1}{\Delta x}
\int_{t^n}^{t^{n+1}}
[F(U)]_{x_{j-1/2}}^{x_{j+1/2}}\,\dd t=0.
\]
No approximation has yet been made.
\sourcefig{page_42_fig_01.png}{0.70}

At an interface, the two boundary fluxes are
\[
F(U)=
\begin{bmatrix}
\rho u\\
\rho u^2+p\\
\rho h u
\end{bmatrix}.
\]
If \(U\) is known at these spots, the integral can be computed. Solve a
Riemann problem at the interface: depending on where the \(t\)-axis lies, the
state is known.

\sourcefig{page_42_fig_02.png}{0.72}

This state is constant only if no other waves are involved. Since
\[
\Delta t^n=\alpha\frac{\Delta x}{\lambda_{j,\max}},
\]
this is guaranteed, and the two interface fluxes are constant. For each cell,
two Riemann problems must be solved, and
\[
[F(U)]_{x_{j-1/2}}^{x_{j+1/2}}
=F_{j+1/2}^{\mathrm{RP}}-F_{j-1/2}^{\mathrm{RP}}.
\]
Therefore
\[
\overline U_j^{\,n+1}-\overline U_j^{\,n}
+\frac{\Delta t^n}{\Delta x}
\left(F_{j+1/2}^{\mathrm{RP}}-F_{j-1/2}^{\mathrm{RP}}\right)=0.
\]

The Godunov scheme is
\[
\boxed{
\overline U_j^{\,n+1}
=\overline U_j^{\,n}
-\frac{\Delta t^n}{\Delta x}
\left(F_{j+1/2}^{\mathrm{RP}}-F_{j-1/2}^{\mathrm{RP}}\right).}
\]
It is stable and robust because it can be intrinsically proved that
\(s\nearrow\), even if it should remain constant; hence it gives a physical
answer. It also ``emulates'' the real behavior of the flow thanks to a
``dissipation'' process given by the fact that it computes small jumps as
almost smooth variations.

% SOURCE PAGE 43
\section*{Boundary Conditions}
\textit{marted\`i 14 gennaio 2025, 20:07}

The assignment of proper boundary conditions is crucial in the resolution of
any kind of flow, homentropic, isentropic, or not isentropic. In gas dynamics
these conditions often depend also on the flow itself; hence they must always
be checked when there is an interaction with the flow. The following are
examples of common boundary conditions for both wave form and matrix form.
Boundary conditions are related to the nature of the problem.

\subsection*{Piston and Wall}
The non-slip and impermeability conditions are
\[
u=v_p\qquad\text{(piston)},\qquad
u=0\qquad\text{(wall)}.
\]

\subsection*{Open End}
\sourcefig{page_43_fig_01.png}{0.64}

The number of boundary conditions depends on whether the point is to the left
of \(A\) or to the right of \(B\), i.e. on the number of characteristic lines
coming from outside.

\sourcefig{page_43_fig_02.png}{0.96}
\sourcefig{page_43_fig_03.png}{0.94}

The source sketches the cases \(u>0,M<1\) (two conditions), \(u=0,M=1\)
(two conditions), \(u>0,M>1\) (three conditions), and a subsonic case with
one condition, and so on.

The boundary condition itself also depends on the nature of the outside.

% SOURCE PAGE 44
\sourcefig{page_44_fig_01.png}{0.96}

\subsection*{Vessel condition}
For a large tank,
\[
p_0=\mathrm{const},\qquad T_0=\mathrm{const},\qquad u=0.
\]
If the flow exits the tube and enters the vessel, the energy is known and
\(p_{\mathrm{exit}}=p_0\).  If flow comes from the vessel, part of the energy
comes from the vessel and the condition is more complex.

If
\[
\left|\frac{\Delta\rho}{\rho_0}\right|=\frac{A L}{V}\ll1,
\qquad
\left|\frac{\Delta p}{p_0}\right|=\gamma\frac{A L}{V}\ll1,
\]
where \(A\) is the tube cross section, \(L\) its length, and \(V\) the tank
volume,
the source uses
\[
p_0=p\left(1+\delta M^2\right)^{\gamma/(\gamma-1)},
\qquad
s_0=s.
\]
Here \(p\) and \(M\) are those of the flow; the other required condition is
\(s_0=s\).

\sourcefig{page_44_fig_02.png}{0.96}

In Example 1, the left vessel has \(p_{01},T_{01},u_{01}=0\), the right vessel
has \(p_{04},T_{04},u_4=0\), and \(p_{01}>p_{04}\), giving \(u>0\). At the
left boundary the flow is coming from the tank; at the right boundary the
flow is going into the tank.

\sourcefig{page_44_fig_03.png}{0.96}

% SOURCE PAGE 45
\sourcefig{page_45_fig_01.png}{0.96}

For state 5, the source writes
\[
p_{01}=p_5(1+\delta M_5^2)^{\gamma/(\gamma-1)},
\qquad s_{01}=s_3,
\]
and
\[
a_5+\delta u_5=a_2+\delta u_2,
\qquad
\frac{p_5}{p_2}
=\left(\frac{a_5}{a_2}\right)^{2\gamma/(\gamma-1)},
\qquad
p_{01}=p_5(1+\delta M_5^2)^{\gamma/(\gamma-1)}.
\]
For state 6,
\[
a_6+\delta u_6=a_3+\delta u_3,
\qquad
\frac{p_6}{p_3}
=\left(\frac{a_6}{a_3}\right)^{2\gamma/(\gamma-1)},
\qquad
p_6=p_{04}.
\]

\subsection*{Example 2}
\sourcefig{page_45_fig_02.png}{0.96}

The initial vessel state is \(p_1,T_1,u_1=0\); the reservoir state is
\(p_0,T_0,s_0,u=0\), with \(p_1>p_0\).

\sourcefig{page_45_fig_03.png}{0.66}

For the first wave, which is in the \(u+a\) family,
\[
a_1+\delta u_1^0=a_2+\delta u_2,
\qquad
\frac{p_2}{p_1}
=\left(\frac{a_2}{a_1}\right)^{2\gamma/(\gamma-1)},
\qquad
p_2=p_0
\]
because the flow goes out.  The wall then imposes \(u_3=0\):
\[
a_2-\delta u_2=a_3-\delta u_3,
\qquad
u_3=0,
\qquad
\frac{p_3}{p_2}
=\left(\frac{a_3}{a_2}\right)^{2\gamma/(\gamma-1)}.
\]

\sourcefig{page_45_fig_04.png}{0.66}

The reflected wave reaches the vessel boundary. It changes nature, so \(u\)
is turned counterclockwise and the flow goes in: the vessel boundary
condition changes.

The next state begins with the reflected-wave invariant
\[
a_3-\delta u_3=a_4-\delta u_4.
\]

% SOURCE PAGE 46
The remaining relations are
\[
\frac{p_4}{p_3}
=\left(\frac{a_4}{a_3}\right)^{2\gamma/(\gamma-1)},
\qquad
p_0=p_4(1+\delta M_4^2)^{\gamma/(\gamma-1)}.
\]
The source continues \lq\lq and so on...\rq\rq.

\subsection*{Matrix Form}
The wave-splitting form is
\[
U_t+R\Lambda^+L\,U_x\big|^+
+R\Lambda^-L\,U_x\big|^-=0.
\]
This formulation allows discontinuities in the gradient of the initial
conditions and separates waves coming from the left from waves coming from
the right.  Depending on the boundary, some information is known from the
inside while other information comes from outside.
At a left boundary, for the incoming waves \(\lambda_3=u+a\) and
\(\lambda_2=u\), the source identifies the wave variables
\[
p_x+\rho a u_x\qquad\text{and}\qquad
\rho_x-\frac{1}{a^2}p_x.
\]
These wave variables are added to the unknowns \(p,u,\rho\), and the
known and unknown terms can be separated as
\[
U_t+R\Lambda^+L\,U_x\big|^+
=-R\Lambda^-L\,U_x\big|^-.
\]
The component equations are
\[
\begin{aligned}
\rho_t+u\left(\rho_x-\frac{1}{a^2}p_x\right)
+\frac{u+a}{2a^2}(p_x+\rho a u_x)
&=-\frac{u-a}{2a^2}(p_x-\rho a u_x),\\
u_t+\frac{u+a}{2\rho a}(p_x+\rho a u_x)
&=\frac{u-a}{2\rho a}(p_x-\rho a u_x),\\
p_t+\frac{u+a}{2}(p_x+\rho a u_x)
&=-\frac{u-a}{2}(p_x-\rho a u_x).
\end{aligned}
\]
The source notes that this gives five unknowns and three equations, so two
equations are still needed.  These boundary equations must be written in
differential form and cannot be assigned directly to \(p,u,\rho\) or to the
wave variables, because they also depend on the interior state.
Two examples are presented in lesson 12/18.

\subsection*{Godunov Scheme}
We have seen that, practically speaking, computing the solution with the
Godunov method requires solving Riemann problems at cell interfaces. To
determine the right condition at the boundary, ghost cells are used. These
are unphysical cells placed outside the domain and purposefully created to
impose the boundary condition.

% SOURCE PAGE 47
\sourcefig{page_47_fig_01.png}{0.72}

\textcolor{green}{\textbf{Example: wall condition}}\quad u=0~\mathrm{m/s}.
Since the solution of the Riemann problem at the boundary must have
\(u_0^0=0~\mathrm{m/s}\), assign the ghost-cell velocity
\[
\overline u_g^0:=-\overline u_0^0
\]
so that the imaginary resultant flow has zero velocity.

% SOURCE PAGE 48
\section*{Quasi 1-D}
\textit{giovedì 16 gennaio 2025, 17:18}

This is a simplified axially symmetric model.  The axial velocity \(u\) is
much larger than \(v,w\), and the area varies in \(x\); variables are
cross-sectional averages.

\subsection*{Isentropic Model (No shocks, differential form)}
\[
\begin{aligned}
(\rho A)_t+(\rho u A)_x&=0
&&\text{mass},\\
(\rho u A)_t+\bigl[((\rho u)u+p)A\bigr]_x-pA_x&=0
&&\text{momentum},\\
(\rho e A)_t+(\rho e u A+p u A)_x&=0
&&\text{energy}.
\end{aligned}
\]
In the stationary case, considering the mass equation and the total pressure
and speed of sound, one obtains the area law
\[
\frac{A_2}{A_1}
=\frac{M_1}{M_2}
\left(\frac{1+\delta M_1^2}{1+\delta M_2^2}\right)^{
\frac{\gamma+1}{2(\gamma-1)}}.
\]

\subsection*{Non Isentropic case}
The formulation is complex and it has not been covered this year, but some
general facts have been clarified:
\begin{enumerate}
  \item It is enough to know the initial condition (i.e. states 2 and
        \(1=\mathrm{amb}\)) to know the strength of the shock.
  \item When \(A\downarrow\), the shock becomes stronger; on the contrary, if
        \(A\uparrow\), it weakens.  This is because of the different
        space where the finite energy is distributed.
  \item Since \(W\) changes with \(A\), the entropy jump does the same and an
        entropy layer is formed.
  \item If the initial shock is strong enough, sonic conditions are reached at
        the throat.  If the pressure ratio is enough, a supersonic flow
        develops past the throat; what is behind the throat does not receive
        any information about what is past the throat.
\end{enumerate}

\sourcefig{page_48_fig_01.png}{0.66}

% SOURCE PAGE 49
\subsection*{Nozzle Start Up (simplified)}
Consider an abrupt area variation.  Assume the shock is strong enough that
\(M_3\geq 1\) at the throat.

\sourcefig{page_49_fig_01.png}{0.68}
\sourcefig{page_49_fig_02.png}{0.68}

When it crosses the throat, the flow adapts to the change of \(A\) by
generating waves.  Since \(M_3\geq 1\), the expected pattern is that shown
above.  To understand what the third wave is, use the \(p\)-\(u\) plane; it is
found to be another shock.

\sourcefig{page_49_fig_03.png}{0.68}

In the \(p\)-\(u\) plane, the area law computes state 5, which corresponds to
the instant after the change of \(A\).  This law can be used because the flow
is supersonic and nothing can come back from state 1.  The resulting
phenomenon is shown below.

\sourcefig{page_49_fig_04.png}{0.68}
\sourcefig{page_49_fig_05.png}{0.72}

% SOURCE PAGE 50
\section*{2-D Steady Flows}
\textit{giovedì 16 gennaio 2025, 19:16}

\subsection*{Balance Law and characteristic theory (HOMENTHALPIC and HOMENTROPIC)}
This theory can be applied to three-dimensional problems with axial symmetry.
The general divergence form is
\[
U_t+F(U)_x+G(U)_y=0.
\]
With
\[
M=\frac{\sqrt{u^2+v^2}}{a},
\qquad
U=\begin{bmatrix}\rho\\\rho u\\\rho v\\\rho e\end{bmatrix},
\]
\[
F(U)=
\begin{bmatrix}
\rho u\\
(\rho u)u+p\\
(\rho u)v\\
\rho e u+p u
\end{bmatrix},
\qquad
G(U)=
\begin{bmatrix}
\rho v\\
(\rho u)v\\
(\rho v)v+p\\
\rho e v+p v
\end{bmatrix}.
\]
For a stationary flow, \(F_x+G_y=0\).  Since time is not important,
isentropicity and isenthalpicity are granted by definition; their spatial
distribution must instead be assumed.  No discontinuities are considered for
now.

\sourcefig{page_50_fig_01.png}{0.50}

Writing
\[
A^xU_x+A^yU_y=0,
\]
if \(\theta\ll 1\), it can be proved that the system becomes hyperbolic with
\(A=(A^x)^{-1}A^y\):
\[
U_x+A U_y=0.
\]
This is mathematically very similar to the one-dimensional unsteady system.
Hence waves, directions, and transported variables can be found if the system
is reducible (its eigenvalues are real and nonzero).  A straightforward
approach would diagonalize \(A\), but proceed instead from the expanded
primitive equations:
\[
\begin{aligned}
u\rho_x+v\rho_y+\rho(u_x+v_y)&=0, \tag{I}\\
\rho u u_x+\rho v u_y+p_x&=0, \tag{II}\\
\rho u v_x+\rho v v_y+p_y&=0, \tag{III}\\
u p_x+v p_y+\rho a^2(u_x+v_y)&=0. \tag{IV}
\end{aligned}
\]
Subtracting \(a^2\) times (I) from (IV),
\[
u(p_x-a^2\rho_x)+v(p_y-a^2\rho_y)
+(\rho a^2-\rho a^2)(u_x+v_y)=0
\qquad\Longrightarrow\qquad
u s_x+v s_y=0.
\]
Thus entropy is constant along streamlines,
\(\vect V\mathbin{\cdot}\nabla s=0\).

% SOURCE PAGE 51
Consider
\[
h:=\frac{p}{\gamma-1}+\rho\frac{|\vect V|^2}{2}+p,
\]
including internal, kinetic, and pressure-force contributions.  Manipulating
(IV) gives
\((\vect V\mathbin{\cdot}\nabla)h=0\), so enthalpy is also constant along
streamlines (constant in space, i.e. uniform along streamlines).  The first two
wave variables are \(h\) and \(s\), transported by the streamlines (velocities).

The other two wave variables come from the eigenproblem.  With flow angle
\(\theta\), Mach angle \(\alpha\), and \(M>1\), the eigenvalues are
\[
\lambda_1=\tan(\theta-\alpha),\qquad
\lambda_4=\tan(\theta+\alpha),
\qquad
\alpha=\sin^{-1}\frac{1}{M}.
\]
The corresponding differential relations are
\[
-v\,\dd u+u\,\dd v-\frac{1}{\rho}\sqrt{M^2-1}\,\dd p=0,
\]
\[
-v\,\dd u+u\,v+\frac{1}{\rho}\sqrt{M^2-1}\,\dd p=0.
\]

\sourcefig{page_51_fig_01.png}{0.48}

At \(M=1\) the two equations coincide (the degenerate case).  They can be
integrated---that is, conserved variables can be defined and the system is
reducible---only if the flow is homenthalpic, homentropic, and \(M\neq1\).  In
that case,
\[
\theta+\nu(M)=\mathrm{const}
\quad\Longrightarrow\quad \Delta\theta+\Delta\nu=0
\quad\text{along}\quad
\frac{\dd y}{\dd x}=\tan(\theta-\alpha),
\]
\[
\theta-\nu(M)=\mathrm{const}
\quad\Longrightarrow\quad \Delta\theta-\Delta\nu=0
\quad\text{along}\quad
\frac{\dd y}{\dd x}=\tan(\theta+\alpha),
\]
where \(\nu(M)\) is the Prandtl function (p.~135, Franco).

\sourcefig{page_51_fig_02.png}{0.82}

This function links the local Mach number and the deviation of the flow from
the \(x\)-direction.  It is limited to about \(130^\circ\), where
\(M\to\infty\); this can be interpreted as follows.

\sourcefig{page_51_fig_03.png}{0.78}

If \(\theta\to130^\circ\), a void is created (cavitation); hence \(a\to0\).

\subsection*{Czocco Theorem}
This is an important result that links enthalpy, entropy, and vorticity.

% SOURCE PAGE 52
The momentum equation is
\[
(\vect V\mathbin{\cdot}\nabla)\vect V+\frac{1}{\rho}\nabla p=0.
\]
Use the Lagrange identity
\[
(\vect V\mathbin{\cdot}\nabla)\vect V
=\nabla\left(\frac{|\vect V|^2}{2}\right)
+(\nabla\times\vect V)\times\vect V
\]
and
\[
\frac{1}{\rho}\nabla p=\nabla h_i-T\nabla s,
\qquad
h_i=\frac{\gamma}{\gamma-1}\frac{p}{\rho}.
\]
The source obtains
\[
\nabla\left(\frac{|\vect V|^2}{2}+h_i\right)
=\nabla h=T\nabla s-\vect\omega\times\vect V,
\qquad
\vect\omega=\nabla\times\vect V.
\]
Projecting Czocco's equation on the velocity shows that, along streamlines,
the zero-vorticity information is transported.  If the flow is homenthalpic,
\[
\nabla h=T\nabla s-\vect\omega\times\vect V=0
\qquad\Longrightarrow\qquad
T\nabla s=\vect\omega\times\vect V.
\]
The source states that the flow is \lq\lq omentropic\rq\rq if and only if it
is also irrotational; the handwritten word is retained as written.

If entropy gradients (not jumps, since \(s\) is constant) are present,
vorticity is present.

% SOURCE PAGE 53
\section*{Expansion and Compression}
\textit{venerdì 17 gennaio 2025, 09:44}

For steady isentropic flow,
\[
p_0=p(1+\delta M^2)^{\gamma/(\gamma-1)}.
\]
The pressure behind an isentropic wave is found from \(M\), and \(M\) from
\(\nu(M)\) and the invariants
\[
\Delta\theta+\Delta\nu(M)=0
\quad\text{along}\quad
\frac{\dd y}{\dd x}=\theta-\alpha,
\]
\[
\Delta\theta-\Delta\nu(M)=0
\quad\text{along}\quad
\frac{\dd y}{\dd x}=\theta+\alpha,
\qquad
\alpha=\sin^{-1}\frac1M.
\]

\subsection*{Prandtl--Meyer Expansion}
\sourcefig{page_53_fig_01.png}{0.88}

For an abrupt variation of the geometry, the wall turns downward.  Since the
flow is supersonic, both fronts of the generated wave are in the
\(\theta+\alpha\) direction.  The corresponding invariant gives, behind
state 2,
\[
\Delta\theta+\Delta\nu(M)=0,
\qquad \Delta\nu(M)=-\Delta\theta>0,
\qquad M_2>M_1,\qquad p_2<p_1,
\qquad \alpha_2=\sin^{-1}\frac1{M_2}<\alpha_1
\quad\text{(with respect to \(y\))}.
\]
The last pressure result follows from the isentropic pressure relation above.
If the expansion had been smooth, nothing would change in the determination
of state 2; finding the whole \(\vect V\) field would become much harder, but
that is not of interest here.

\subsection*{Compression}
\sourcefig{page_53_fig_02.png}{0.88}

The same construction and formulas apply.  The only difference is
\(\Delta\theta>0\), which leads to the manuscript's
\(\nu(M)<0\), hence \(M_2<M_1\) and \(p_2>p_1\).

% SOURCE PAGE 54
\section*{Oblique Shocks}
\textit{venerdì 17 gennaio 2025, 10:30}

\sourcefig{page_54_fig_01.png}{0.72}

To obtain jump relations in two dimensions, follow the same reasoning as in
one dimension, beginning with the stationary divergence form
\(F(U)_x+G(U)_y=0\).  Since shocks are a plane phenomenon, rotate the
reference frame with the shock through \(\sigma\).  In this frame the
phenomenon is plane, so \(G(U)_\eta=0\), and
\[
\begin{aligned}
(\rho u_n)_\xi&=0,\\
(p+\rho u_n^2)_\xi&=0,\\
(\rho u_n v_t)_\xi&=0,\\
(\rho u_n h)_\xi&=0.
\end{aligned}
\]
The jump relations are
\[
\boxed{
\begin{aligned}
[\rho u_n]&=0,\\
[p+\rho u_n^2]&=0,\\
[\rho u_n v_t]&=0,\\
[\rho u_n h]&=0.
\end{aligned}}
\]
In particular,
\[
\avg{\rho}[u_n]+[\rho]\avg{u_n}=0.
\]
The three possibilities are:
\begin{enumerate}
  \item If \([\rho]=0\) (with \(\rho\neq0\)) and \([u_n]=0\), there is no
        jump of normal velocity or density.  Substitution in the momentum
        relation along \(\xi\) gives \([p]=0\) for every \(\avg{u_n}\): the
        classical contact discontinuity.
  \item If \([\rho]\neq0\), then \(\avg{u_n}=0\) and \([u_n]=0\), hence
        \(u_n=0\).  Momentum along \(\eta\) is satisfied for every
        \([v_t]\): there is no flux and the tangential velocity may jump
        arbitrarily, giving a slip line.  Also \([p]=0\).
  \item The remaining mass relation
\[
\avg{\rho}[u_n]=-[\rho]\avg{u_n}
\]
        and momentum along \(\eta\) give \([v_t]=0\): an oblique shock.
\end{enumerate}

\sourcefig{page_54_fig_02.png}{0.72}

Across an oblique shock, the tangential component of velocity is unchanged,
while the normal component has a jump.

% SOURCE PAGE 55
Across an oblique shock the tangential velocity is unchanged while the normal
velocity jumps:
\[
M_{1n}=\frac{u_n}{a_1}
=\frac{|\vect V_1|\sin\sigma}{a_1}
=M_1\sin\sigma,
\]
\[
M_{1t}=\frac{v_t}{a_1}
=\frac{|\vect V_1|\cos\sigma}{a_1}
=M_1\cos\sigma.
\]
Only the normal Mach number changes.  The normal Rankine--Hugoniot formulas
are
\[
\frac{p_2}{p_1}
=1+\frac{2\gamma}{\gamma+1}(M_{1n}^2-1),
\]
\[
\frac{\rho_2}{\rho_1}
=\frac{u_{n1}}{u_{n2}}
=\frac{\gamma+1}{\gamma-1+2/M_{1n}^2},
\]
\[
M_{2n}^{\,2}
=\frac{\delta+1/M_{1n}^2}{\gamma-\delta/M_{1n}^2}.
\]
The shock angle satisfies
\[
\alpha\leq\sigma\leq\frac{\pi}{2},
\qquad
\sigma=\sin^{-1}\frac{M_n}{M}.
\]
If \(M_n=1\), \(\alpha=\sin^{-1}(1/M)\).  Although \(V_{1t}=V_{2t}\),
\[
M_{1t}=\frac{V_{1t}}{a_1}\neq
\frac{V_{2t}}{a_2}=M_{2t}
\]
because the sound speed changes.

\sourcefig{page_55_fig_01.png}{0.82}

Since
\[
M_{2n}=M_2\sin(\sigma-\theta),
\qquad
M_2=\frac{M_{2n}}{\sin(\sigma-\theta)},
\]
the theta-deflection relation is
\[
\tan\theta
=2\cot\sigma\,
\frac{M_1^2\sin^2\sigma-1}
{(\gamma+1)M_1^2-2M_1^2\sin^2\sigma+2}.
\]
As \(M_1\to\infty\),
\[
\lim\tan\theta
=\frac{2\cot\sigma\sin^2\sigma}
{\gamma+1-2\sin^2\sigma}.
\]
The maximum occurs at approximately
\(\sigma_{\max}=61.8^\circ\), giving
\(\theta_{\max}\simeq45.6^\circ\).

\sourcefig{page_55_fig_02.png}{0.84}

For fixed \(M_1\), with \(\sigma\) dependent on \(\theta\) (shock inclination
as a function of ramp inclination), there can be two solutions.  If
\(\theta>\theta_{\max}\) is imposed, the shock is detached (blunt body).  If
\(\theta<\theta_{\max}\) and \(\sigma<\sigma_{\max}\), the weak solution has
supersonic flow past the shock.  If \(\theta<\theta_{\max}\) and
\(\sigma>\sigma_{\max}\),

% SOURCE PAGE 56
\sourcefig{page_56_fig_01.png}{0.84}

the strong solution has subsonic flow past the shock.

\subsection*{Odograph plane (polar representation)}
For these purposes, use the strong solution inside pipes and the weak solution
in an open environment.  Another ``useful'' representation for discussing the
oblique stationary shock is the polar one.  It is portrayed in the odograph
plane, which requires the concept of critical Mach number.

Define the critical Mach number
\[
M_*:=\frac{u}{a_*},
\]
where \(a_*\) is the sound speed at \(M=1\).  The stagnation relation is
\[
a_0^2=a^2+\delta u^2=\mathrm{const}.
\]
At \(M=1\), \(u_*=a_*\), so
\[
a_0^2=a_*^2+\delta a_*^2,
\qquad
a_*^2=\frac{2}{\gamma+1}a_0^2
=\frac{2}{\gamma+1}(a^2+\delta u^2).
\]
Therefore
\[
\frac{\gamma+1}{2}
=\frac{a^2}{a_*^2}-\delta\frac{u^2}{a_*^2}
=\left(\frac{a}{a_*}\right)^2-\delta M_*^2,
\]
\[
\left(\frac{a}{a_*}\right)^2
=\frac{\gamma+1}{2}+\frac{\gamma-1}{2}M_*^2,
\]
\[
\frac{M_*^2}{M^2}
=\frac{\gamma+1}{2}+\frac{\gamma-1}{2}M_*^2.
\]
Hence
\[
\boxed{M_*^2
=\frac{\gamma+1}{\gamma-1+2/M^2}}.
\]
Useful limits are
\[
M=1\Rightarrow M_*=1,\qquad
M\to\infty\Rightarrow M_*\to\frac{\gamma+1}{\gamma-1},
\qquad
M\to0\Rightarrow M_*\to0.
\]
Across an oblique shock,
\[
M_{1t*}=\frac{V_{1t}}{a_*}
=\frac{V_{2t}}{a_*}=M_{2t*}.
\]
In the odograph plane, with axes \(M_{n*},M_{t*}\), the shock is represented
by
\[
\frac{V_t}{a_*}=M_{t*}
=f\!\left(\frac{u_n}{a_*}\right)=f(M_{n*}),
\]
obtained by combining the Prandtl relation \(M_{1n*}M_{2n*}=1\) with the
normal-shock relation.

% SOURCE PAGE 57
\sourcefig{page_57_fig_01.png}{0.74}

The upper branch of interest has \(M_{n*},M_{t*}>0\).  Fixing \(M_{1*}\)
before the shock gives the curve shown: point \(E\) is the normal shock
(maximum strength), while \(F\) is no shock (zero strength).  For these two
points, \(\overline{OF}\) and \(\overline{OE}\) are \(M_{2*}\) after the
shock.  The curve is the shock polar: the locus of states after every possible
oblique shock for a fixed state before it.  Choosing \(\theta\), considered as
\(\Delta\theta\) from the horizontal, gives three
intersections: \(A\) (strong), \(B\) (weak), and \(C\) (non-admissible bow
shock).  Point \(A\) lies inside the \(M_*=1\) circumference, while \(B\)
does not (\(M_{2*}\geq1\)).

\sourcefig{page_57_fig_02.png}{0.72}

Point \(D\) gives \(\theta_{\max}\) for fixed \(M_{1*}\).  The intersection
between \(M_{1*}\) and the polar is very close to \(\theta_{\max}\), but
slightly less, consistently with the previous graph.  The angle \(\sigma\)
is the angle between \(\overline{OH}\) and the \(M_{n*}\)-axis.  State 2
always lies on the polar and depends on \(M_{1*}\).  If \(\theta=0\), either
states 1 and 2 coincide (no shock) or they are at maximum distance (maximum
jump, normal shock).  If \(\theta\neq0\), state 2 lies as shown; draw the line
joining 1 and 2 and its perpendicular through the origin.  These identify
\(H\).

\sourcefig{page_57_fig_03.png}{0.72}
\sourcefig{page_57_fig_04.png}{0.72}

The source labels \(\overline{OH}=V_t\), \(\overline{2H}=u_{n2}\), and
\(\overline{1H}=u_{n1}\); these lines identify \(H\), and hence \(\sigma\).
Mach circles \(M=2,4,\infty\) are shown.  As \(M_1\) increases, so does
\(M_{1*}\), up to the limits: the strong solution is
\[
\sqrt{\frac{\gamma-1}{\gamma+1}}\simeq0.41,
\qquad\text{and the weak solution is}\qquad
\sqrt{\frac{\gamma+1}{\gamma-1}}\simeq2.45.
\]

% SOURCE PAGE 58
\section*{Busemann Polar and reflected stationary shocks}
The dual \(p\)-\(u\) representation is the \(p\)-\(\theta\) Busemann polar.
It is drawn for a fixed \(M_1\) and is useful for reflection and interaction
phenomena.  Its boundary conditions work just as in the \(p\)-\(u\) plane;
isentropic expansion and compression can also be represented.  Here
\(\theta\) is the deflection angle.  The polar contains normal-shock,
zero-strength, strong, maximum-deflection, sonic, weak, and unphysical
branches.

\sourcefig{page_58_fig_01.png}{0.82}

\subsection*{Regular Reflection}
\sourcefig{page_58_fig_02.png}{0.82}

For regular reflection, assume the first shock is a weak solution, so
\(M_2>1\).  A second
shock of the opposite family is required to make \(\vect V\) parallel to the
wall.  In the \(p\)-\(\theta\) plane the two polars intersect.  Weak
solutions with \(M_3>1\) are typically obtained and are highlighted in the
figure.  Depending on \(\theta^{\mathrm{II}}\) and \(M_2\), some strange
solutions may occur.

\sourcefig{page_58_fig_03.png}{0.76}

% SOURCE PAGE 59
\sourcefig{page_59_fig_01.png}{0.70}

\subsection*{Mach Reflection}
\sourcefig{page_59_fig_02.png}{0.76}

Mach reflection has a Mach stem and a slip line.  It occurs at
\(\theta_{\mathrm I}=\theta_N\), the von Neumann angle, which depends on
\(M_1\).  For \(M_1\geq5\), \(\theta_N\simeq20^\circ\).  At
\(\theta_{\mathrm I}=\theta_N\), the
polars intersect on the pressure axis at the normal-shock solution of the
first polar.  The slip line is generated because \(M_4<1\) and \(M_3>1\).

\sourcefig{page_59_fig_03.png}{0.72}
\sourcefig{page_59_fig_04.png}{0.84}

For \(M_1\geq5\), the second polar can fail to reach the pressure axis while
still intersecting the first curve near
\(\theta_D\simeq30^\circ\), giving a bow shock and a curved slip line.
Experimentally there is a hysteresis cycle when \(M_1\) is fixed and
\(\theta\) is varied: starting from small \(\theta\), regular reflection is
encountered up to \(\theta_D\) and the bow solution; starting from high
\(\theta\), Mach reflection persists until \(\theta_N\), then regular
reflection occurs.

% SOURCE PAGE 60
\sourcefig{page_60_fig_01.png}{0.84}

The \(\theta\)-\(M_1\) plot shows the regular-reflection and
Mach-reflection curves and their regions.  If the reflected-shock polar
intersects neither the pressure axis nor the incident-shock curve, there is
no reflection and a bow shock is perpendicular to the wall.

\sourcefig{page_60_fig_02.png}{0.74}
\sourcefig{page_60_fig_03.png}{0.84}

% SOURCE PAGE 61
\section*{Bow Shock}
\textit{venerdì 17 gennaio 2025, 18:08}

We have seen that the oblique-shock solution exists only if
\(\theta<\theta_{\max}\).  The following is a qualitative description of
what happens.

\sourcefig{page_61_fig_01.png}{0.84}

Around a blunt body with upstream \(M_\infty>1\), the bow shock detaches and is normal on the
axis.  Its inclination varies along the curve, so the strength varies and
creates an entropy gradient and vorticity behind the shock (Crocco theorem).
Far away, the shock strength tends to zero and the shock approaches a Mach
line.  Immediately behind the nose the flow is subsonic and can see the body;
it slows to the still condition, isentropically along the axis line.  At some
point two sonic lines arise because the shock is not strong enough to impose
the transition.  Two small curved shocks form because the recirculation zone
imposes a small but abrupt variation of velocity direction.  Finally an
unsteady wake is present, but it is confined downstream by the supersonic
nature of the flow.  Consequently the behavior depends only on what occurs
before the shock: no information can travel back.

\sourcefig{page_61_fig_02.png}{0.84}

The Busemann rule places the sonic line on the edge of an equivalent ramp
which starts from the intersection of the Mach lines and achieves the same
jump in \(M\).

% SOURCE PAGE 62
\section*{Transonic Theory}
\textit{venerdì 17 gennaio 2025, 19:24}

\subsection*{Few words on blunt body case}
\sourcefig{page_62_fig_01.png}{0.80}

The transonic range for a blunt body is roughly
\(0.45<M_\infty<1\).  For \(M_\infty\simeq0.45\), the sonic condition is
reached at two points while the rest of the field is subsonic.  As
\(M_\infty\) increases, two supersonic regions develop, bounded by two curved
Mach lines; the maximum local \(M\) is at the points shown.  The unsteady wake
propagates throughout the flow and makes the sonic regions oscillate
(buffeting).  This oscillation of the pressure distribution produces an
oscillation of the pressure force, hence fluttering.  In general, transonic
regime means that both supersonic and subsonic regions are present in the
field.

\subsection*{SMALL DISTURBANCE THEORY}
We want to develop a theory capable of analyzing the throat region of a
nozzle.

In the throat region of a nozzle, assume homenthalpic and homentropic flow.
Crocco's theorem then gives irrotational flow.  The Euler equations reduce to
\[
\vect V\mathbin{\cdot}\nabla\left(\frac{|\vect V|^2}{2}\right)
-a^2\nabla\mathbin{\cdot}\vect V=0,
\qquad
\frac{p}{\rho^\gamma}=\mathrm{const},
\qquad
a^2+\delta|\vect V|^2=\mathrm{const}.
\]
Since the field is irrotational, define a potential function by
\(\vect V=\nabla\overline\Phi\).  The momentum equation becomes
\[
(\overline\Phi_x^{\,2}-a^2)\overline\Phi_{xx}
+(\overline\Phi_y-a^2)\overline\Phi_{yy}
+2\overline\Phi_x\overline\Phi_y\overline\Phi_{xy}=0.
\]
Expand about \(M_\infty=1\):
\[
\overline\Phi=V_\infty[x+\varepsilon\varphi(x,y)],
\qquad
V_\infty=\sqrt{u_\infty^2+v_\infty^2}.
\]
Here \(\varphi\) is the perturbation potential and \(\varepsilon\) is a
geometry-, Mach-number-, and \(\gamma\)-dependent expansion parameter.

% SOURCE PAGE 63
If \(M_\infty\to1\), then \(\varepsilon\to0\) and \(\varepsilon\ll1\).
Rescale
\[
\xi=\frac{x}{l},\qquad
\eta=\frac{\lambda y}{l},
\]
where \(l\) is the horizontal length and \(\lambda\ll1\).  The source notes
that when \(M_\infty\simeq1\), characteristic lines are perpendicular to the
streamlines and need a larger transverse scale.

\sourcefig{page_63_fig_01.png}{0.96}

For a geometry \(y=s f(x)\), with thickness scale \(s\), the source takes
\(\varepsilon=s/\lambda\) and writes
\[
\Phi=V_\infty\left[x+\frac{s}{\lambda}\varphi(x,y)\right].
\]
After substitution and neglect of higher-order terms, the transonic small
disturbance equation is copied as
\[
\left[
\frac{1-M_\infty^2}{\lambda^2}
-\frac{(\gamma+1)\varepsilon M_\infty^2}{\lambda^2}\varphi_\xi
\right]\varphi_{\xi\xi}
+\left[1-(\gamma-1)\varepsilon M_\infty^2\varphi_\xi\right]\varphi_{\eta\eta}
-2M_\infty^2\varepsilon\varphi_\eta\varphi_{\xi\eta}=0.
\]

The sonic scaling is
\[
\frac{\varepsilon M_\infty^2}{\lambda^2}=1,
\qquad
\frac{sM_\infty^2}{\lambda^3}=1.
\]
In the reduced form,
\[
\left[\frac{1-M_\infty^2}{\lambda^2}-(\gamma+1)\varphi_\xi\right]
\varphi_{\xi\xi}+\varphi_{\eta\eta}=0.
\]
Writing \(k\) for the reference state gives
\[
\boxed{[k-(\gamma+1)\varphi_\xi]\varphi_{\xi\xi}
+\varphi_{\eta\eta}=0.}
\]
At the sonic condition \(k=0\), and no singularities are obtained.

% SOURCE PAGE 64
At \(M_\infty=1\), \(k=0\), and
\[
-(\gamma+1)\varphi_\xi\varphi_{\xi\xi}
+\varphi_{\eta\eta}=0
\qquad\Longrightarrow\qquad
(\gamma+1)u u_\xi-v_\eta=0.
\]
Irrotationality gives
\[
v_\xi-u_\eta=0.
\]
We started from \(x\)-\(y\) coordinates and obtained \(u,v\) in the new
coordinates \(\xi\)-\(\eta\), which are linked to \(\varepsilon\).  From now
on consider \(u<0\) subsonic, \(u=0\) sonic, and \(u>0\) supersonic.

The source writes this as a first-order system with
\[
\begin{bmatrix}u\\v\end{bmatrix}_\xi
+A\begin{bmatrix}u\\v\end{bmatrix}_\eta=0,
\qquad
A=
\begin{bmatrix}
0&-\dfrac{1}{(\gamma+1)u}\\
-1&0
\end{bmatrix},
\]
whose eigenvalues are
\[
\lambda=\pm\frac{1}{\sqrt{(\gamma+1)u}}.
\]

For \(u>0\) the system is hyperbolic; at \(u=0\), \(\lambda\to\infty\) and
the characteristics are perpendicular to the streamlines; for \(u<0\) the
system is elliptic.

Therefore this system is consistent with previous descriptions
(``l'utilità di fare sta roba così un giorno forse mi sarà chiara dio cristo'').

Finally, the source asks whether this method can account for discontinuities.
It answers yes, although these are not proper shocks because the flow was
assumed homentropic.  The idea is to return to a divergence form that can be
integrated to obtain jump relations.  Considering a direction \(w\) of the
discontinuity gives
\[
w=\pm\frac{1}{\sqrt{(\gamma+1)\avg{u}}}
\qquad\text{(totalmente inutili/unphysical)}.
\]

\subsection*{Sauer--Meyer Solution for Nozzles}
\sourcefig{page_64_fig_01.png}{0.86}

Consider a nozzle in the propulsive regime, with a symmetric throat region
defined by two arcs of circumference of radius \(R\gg1\).  Use the following
expansion of the perturbation potential:
\[
\varphi(\xi,\eta)=f_0(\xi)+f_1(\xi)\eta+f_2(\xi)\eta^2
+f_3(\xi)\eta^3+f_4(\xi)\eta^4.
\]

% SOURCE PAGE 65
The polynomial displayed at the end of the preceding source page is repeated
at the top of this page.  The idea is to consider placement of the reference
frame as in the figure for symmetry reasons, while keeping the position of the
origin unknown.
Write
the equations for \(u\) and \(v\) in this reference frame, considering the
potential expanded as presented and also a linear distribution of \(u\) along
the \(\xi\)-axis:
\[
u(\xi,0)=d\xi.
\]

Replacing everything in the TSD equation and skipping three pages of passages,
the source obtains
\[
u=d\xi+\frac{d^2}{2}\eta^2,
\qquad
v=d^2\xi\eta+\frac{d^3}{6}\eta^3,
\]
with a red asterisk and the word ``below'' beside the first formula.

These expressions allow computation of the sonic-line curve, the throat-line
curve, and the horizontal position of the origin.

The boundary condition is \(\vect V\) parallel to the nozzle wall, giving
\[
d=\frac{1}{\sqrt R}.
\]
The throat condition \(v=0\) gives
\[
\xi_t=-\frac{d}{6}\eta^2,
\]
a parabola touching the throat, with vertex on the symmetry axis
\(\eta=1\).  The sonic condition \(u=0\) gives
\[
\xi_s=-\frac{d}{2}\eta^2.
\]
For \(R\gg1\), \(\eta\simeq1\), so the parabola starts at approximately
\(-d/2\).

The source marks a 2024 exercise by ``Gzossi'':
\[
d=\frac{1}{\sqrt{(\gamma+1)R}},
\]
\[
\begin{aligned}
u&=d\xi+\frac{(\gamma+1)d^2}{2}\eta^2,\\
v&=(\gamma+1)d^2\xi\eta
+\frac{(\gamma+1)d^3}{6}\eta^3,
\end{aligned}
\qquad
M=\sqrt{(1+u)^2+v^2}.
\]
The source highlights the \(\gamma+1\) factor and notes that Franco omitted
it; it assumes the exercise version is correct even if Favini wrote a
different form.

% ============================================================
% TRANSCRIPTION REPORT
% ============================================================
% Source pages processed: 1--65 (65 pages)
% Pages with unresolved ambiguities: none after direct visual audit
% Pages containing [ILLEGIBLE] markers: none
% Pages marked as blank or continuation: 21
% Pages with extracted figures: 3, 7--15, 17, 19--20, 22--24, 26--35, 37, 41--45, 47--49, 50--55, 57--64
% Total extracted figures: 128
% Figures requiring manual crop review: all 128 crops were reviewed directly against their source pages; source coordinates are recorded in figures/manifest.json
% Figures intentionally not extracted: pages without a distinct source drawing or with only text/formulas
% Other fidelity issues requiring manual review: none after two complete visual passes
% Second-pass audit: complete, source pages 1--65 re-verified against this final version

\end{document}
